%%%%%%%%%%%%%%%%%%%%%%%%%%%%%%%%%%%%%%%%%%%%%%%%%%%%%%%%%%%%
%% Abstract Algebra II — Fields and Galois Theory
%% Undergraduate L3
%% Part 1: Preamble through end of Chapter 10
%%%%%%%%%%%%%%%%%%%%%%%%%%%%%%%%%%%%%%%%%%%%%%%%%%%%%%%%%%%%
\documentclass[12pt,a4paper,openany]{book}

%% ── Encoding and fonts ──────────────────────────────────
\usepackage[utf8]{inputenc}
\usepackage[T1]{fontenc}

%% ── Mathematics ─────────────────────────────────────────
\usepackage{amsmath}
\usepackage{amssymb}
\usepackage{amsthm}
\usepackage{mathtools}
\usepackage{mathrsfs}

%% ── Graphics and diagrams ───────────────────────────────
\usepackage{tikz}
\usepackage{tikz-cd}
\usetikzlibrary{positioning, arrows.meta, calc}

%% ── Colored boxes ───────────────────────────────────────
\usepackage[most]{tcolorbox}
\tcbuselibrary{theorems,skins,breakable}

%% ── Page layout ─────────────────────────────────────────
\usepackage[margin=2.5cm]{geometry}
\usepackage{fancyhdr}
\pagestyle{fancy}
\fancyhf{}
\fancyhead[LE]{\textsc{\nouppercase{\leftmark}}}
\fancyhead[RO]{\textsc{\nouppercase{\rightmark}}}
\fancyfoot[C]{\thepage}
\renewcommand{\headrulewidth}{0.4pt}

%% ── Lists, indexing, links, tables, colour ──────────────
\usepackage{enumitem}
\usepackage{makeidx}
\makeindex
\usepackage[colorlinks=true,linkcolor=blue!70!black,
            citecolor=green!50!black,urlcolor=blue!60!black]{hyperref}
\usepackage{booktabs}
\usepackage{xcolor}
\usepackage{caption}

%% ══════════════════════════════════════════════════════════
%% Theorem-like environments (amsthm)
%% ══════════════════════════════════════════════════════════
\theoremstyle{definition}
\newtheorem{definition}{Definition}[chapter]
\newtheorem{theorem}{Theorem}[chapter]
\newtheorem{proposition}{Proposition}[chapter]
\newtheorem{lemma}{Lemma}[chapter]
\newtheorem{corollary}{Corollary}[chapter]

\theoremstyle{remark}
\newtheorem{remark}{Remark}[chapter]
\newtheorem{example}{Example}[chapter]
\newtheorem{exercise}{Exercise}[chapter]

%% ══════════════════════════════════════════════════════════
%% tcolorbox styling for selected environments
%% ══════════════════════════════════════════════════════════
\tcolorboxenvironment{definition}{
  enhanced,
  breakable,
  colback=blue!5,
  colframe=blue!70!black,
  boxrule=0.8pt,
  left=6pt, right=6pt, top=6pt, bottom=6pt,
  sharp corners,
  before skip=10pt, after skip=10pt,
  fonttitle=\bfseries
}

\tcolorboxenvironment{theorem}{
  enhanced,
  breakable,
  colback=red!5,
  colframe=red!70!black,
  boxrule=0.8pt,
  left=6pt, right=6pt, top=6pt, bottom=6pt,
  sharp corners,
  before skip=10pt, after skip=10pt,
  fonttitle=\bfseries
}

\tcolorboxenvironment{proposition}{
  enhanced,
  breakable,
  colback=green!5,
  colframe=green!60!black,
  boxrule=0.8pt,
  left=6pt, right=6pt, top=6pt, bottom=6pt,
  sharp corners,
  before skip=10pt, after skip=10pt,
  fonttitle=\bfseries
}

%% ══════════════════════════════════════════════════════════
%% Custom commands
%% ══════════════════════════════════════════════════════════
\newcommand{\N}{\mathbb{N}}
\newcommand{\Z}{\mathbb{Z}}
\newcommand{\Q}{\mathbb{Q}}
\newcommand{\R}{\mathbb{R}}
\newcommand{\C}{\mathbb{C}}
\newcommand{\F}{\mathbb{F}}
\DeclareMathOperator{\Gal}{Gal}
\DeclareMathOperator{\Aut}{Aut}
\DeclareMathOperator{\Hom}{Hom}
\DeclareMathOperator{\Ker}{Ker}
\DeclareMathOperator{\im}{im}
\DeclareMathOperator{\id}{id}
\DeclareMathOperator{\chr}{char}
\renewcommand{\deg}{\operatorname{deg}}
\DeclareMathOperator{\irr}{irr}
\DeclareMathOperator{\Fix}{Fix}
\newcommand{\longmapsfrom}{\longleftarrow\mkern-2mu\mathrel{|}}

%% ══════════════════════════════════════════════════════════
%% Title data
%% ══════════════════════════════════════════════════════════
%% ══════════════════════════════════════════════════════════
\begin{document}
\begin{titlepage}
\begin{tikzpicture}[remember picture, overlay]
  \fill[left color=blue!15!white, right color=blue!5!white]
    (current page.south west) rectangle (current page.north east);
  \fill[color=orange!70!yellow]
    ([yshift=-2cm]current page.north west) rectangle ([yshift=-2.3cm]current page.north east);
  \fill[color=blue!80!black, rounded corners=8pt]
    ([xshift=3cm, yshift=-4cm]current page.north west) rectangle
    ([xshift=-3cm, yshift=-10cm]current page.north east);
  \node[anchor=center, text=white, font=\Huge\bfseries]
    at ([yshift=-6cm]current page.north) {Abstract Algebra II};
  \node[anchor=center, text=orange!80!yellow, font=\Large]
    at ([yshift=-7.5cm]current page.north) {Lecture Notes};
  \node[anchor=center, text=white!80, font=\large]
    at ([yshift=-8.8cm]current page.north) {Licence L3 --- 2025--2026};
  \node[anchor=center, text=blue!80!black, font=\Large\itshape]
    at ([yshift=-12cm]current page.north) {Ya\'e Ulrich Gaba};
  \draw[orange!70!yellow, line width=2pt]
    ([xshift=5cm, yshift=-13.5cm]current page.north west) --
    ([xshift=-5cm, yshift=-13.5cm]current page.north east);
  \node[anchor=center, text width=12cm, align=center, text=blue!60!black, font=\itshape]
    at ([yshift=-16cm]current page.north) {``Since Galois, algebra is no longer the art of solving equations, but the study of structures.''\\[4pt]--- Nicolas Bourbaki};
  \node[anchor=south, text=gray, font=\small]
    at ([yshift=1.5cm]current page.south) {\today};
\end{tikzpicture}
\end{titlepage}
%% ── Table of Contents ───────────────────────────────────
\tableofcontents
\newpage

%% ══════════════════════════════════════════════════════════
%% PREFACE
%% ══════════════════════════════════════════════════════════
\chapter*{Preface}
\addcontentsline{toc}{chapter}{Preface}
\markboth{Preface}{Preface}

\index{Galois, Evariste@Galois, \'Evariste}

On the morning of 30 May 1832, a twenty-year-old mathematician named
\'Evariste Galois\index{Galois, Evariste@Galois, \'Evariste} was
fatally wounded in a duel. The night before, fully expecting to die, he
had spent hours feverishly writing letters and manuscripts, scribbling
in the margins \emph{``I have no time''} as he raced to set down the
ideas that had consumed him for years. Those hastily written pages would
eventually revolutionise mathematics.

Galois had discovered a profound connection between two seemingly
disparate areas of algebra: the theory of polynomial equations and the
theory of groups. He showed that to every polynomial equation one can
attach a group --- now called the \emph{Galois group}\index{Galois
group} --- and that the solvability of the equation by radicals is
completely determined by the structure of this group. In doing so, he
not only settled the centuries-old problem of solving polynomial
equations of degree five and higher, but he also laid the foundations for
what we now call \emph{abstract algebra}.

This course, \emph{Abstract Algebra~II}, is devoted to the theory of
fields and the Galois correspondence. It is a continuation of a first
course in abstract algebra (groups and rings) and is designed for
third-year undergraduate students. Our journey will proceed as follows:

\begin{itemize}[leftmargin=2em]
  \item \textbf{Chapters 8--10.} We develop the theory of field
    extensions: algebraic and transcendental elements, splitting fields,
    separability, algebraic closure, and the structure of finite fields.
  \item \textbf{Chapters 11--13.} We establish the Galois
    correspondence and apply it: normal extensions, the fundamental
    theorem, and computational techniques for determining Galois groups.
  \item \textbf{Chapter 14.} We reach the summit: the proof that the
    general polynomial of degree~$\ge 5$ is not solvable by radicals.
\end{itemize}

The prerequisites are a solid understanding of groups, rings, and ideals
at the level of a standard first course. We make free use of quotient
rings, polynomial rings, and the theory of principal ideal domains.

The reader is encouraged to attempt every exercise; it is only through
sustained effort that the beauty and power of Galois theory reveals
itself. Many of the exercises develop important examples or complete
technical details left to the reader in the text.

\bigskip
\hfill\textit{The Authors}

%% ══════════════════════════════════════════════════════════
%% NOTATION
%% ══════════════════════════════════════════════════════════
\chapter*{Notation}
\addcontentsline{toc}{chapter}{Notation}
\markboth{Notation}{Notation}

\begin{center}
\renewcommand{\arraystretch}{1.4}
\begin{tabular}{@{} l p{10cm} @{}}
\toprule
\textbf{Symbol} & \textbf{Meaning} \\
\midrule
$\N, \Z, \Q, \R, \C$ & Natural numbers, integers, rationals, reals,
  complex numbers \\
$\F_q$ & Finite field with $q$ elements \\
$K, L, M, F$ & Fields \\
$L/K$ & Field extension: $L$ is an extension of $K$ \\
$[L:K]$ & Degree of the extension $L/K$ \\
$K(\alpha), K(\alpha_1,\dots,\alpha_n)$ & Field generated by
  $\alpha$ (resp.\ $\alpha_1,\dots,\alpha_n$) over $K$ \\
$K[\alpha]$ & Ring generated by $\alpha$ over $K$ \\
$K[X]$ & Polynomial ring over $K$ in one indeterminate \\
$\irr(\alpha,K)$ & Minimal polynomial of $\alpha$ over $K$ \\
$\deg f$ & Degree of a polynomial $f$ \\
$\chr(K)$ & Characteristic of a field $K$ \\
$\Gal(L/K)$ & Galois group of the extension $L/K$ \\
$\Aut(L)$ & Automorphism group of $L$ \\
$\Aut_K(L)$ & Group of $K$-automorphisms of $L$ \\
$\Hom_K(L,M)$ & Set of $K$-algebra homomorphisms from $L$ to $M$ \\
$\Fix(H)$ & Fixed field of a subgroup $H \le \Aut(L)$ \\
$\id$ & Identity map \\
$f'$ & Formal derivative of the polynomial $f$ \\
$\langle S \rangle$ & Subgroup (or ideal, or subfield) generated by
  $S$ \\
$|G|$, $[G:H]$ & Order of a group $G$, index of $H$ in $G$ \\
$G \cong H$ & $G$ is isomorphic to $H$ \\
$G \rtimes H$ & Semidirect product \\
$S_n, A_n, D_n, C_n$ & Symmetric, alternating, dihedral, cyclic groups \\
\bottomrule
\end{tabular}
\end{center}

%% ══════════════════════════════════════════════════════════
%%
%%   PART III: FIELDS AND GALOIS THEORY
%%
%% ══════════════════════════════════════════════════════════
\part{Fields and Extensions}

%% ══════════════════════════════════════════════════════════
%% CHAPTER 8: FIELDS — ALGEBRAIC AND TRANSCENDENTAL EXTENSIONS
%% ══════════════════════════════════════════════════════════
\chapter{Fields --- Algebraic and Transcendental Extensions}
\label{ch:field-extensions}
\index{field extension}

\section{Field extensions and degree}
\label{sec:field-ext-degree}

\begin{definition}[Field extension]
\label{def:field-extension}
\index{field extension!definition}
A \textbf{field extension}\index{field extension} is a pair of fields
$K \subseteq L$ (equivalently, an injective ring homomorphism
$\iota\colon K \hookrightarrow L$). We write $L/K$ and call $K$ the
\textbf{base field}\index{base field} and $L$ the \textbf{extension
field}\index{extension field}.
\end{definition}

Any field extension $L/K$ makes $L$ into a $K$-vector space via the
multiplication of $L$.

\begin{definition}[Degree of an extension]
\label{def:degree-extension}
\index{degree!of an extension}
The \textbf{degree} of the extension $L/K$, denoted $[L:K]$, is the
dimension of $L$ as a $K$-vector space:
\[
  [L:K] \;=\; \dim_K L.
\]
If $[L:K] < \infty$, we say $L/K$ is a \textbf{finite
extension}\index{finite extension}. Otherwise it is an
\textbf{infinite extension}\index{infinite extension}.
\end{definition}

\begin{example}[Basic examples]
\label{ex:basic-extensions}
\begin{enumerate}[label=(\roman*)]
  \item $\C/\R$ has degree $[\C:\R] = 2$, with basis $\{1, i\}$.
  \item $\R/\Q$ is an infinite extension (in fact $\R$ is uncountable
    while $\Q$ is countable).
  \item For any field $K$, the extension $K/K$ has degree $1$.
  \item $\Q(\sqrt{2})/\Q$ has degree $2$, with basis $\{1, \sqrt{2}\}$.
\end{enumerate}
\end{example}

\begin{theorem}[Tower Law]
\label{thm:tower-law}
\index{tower law}
Let $K \subseteq L \subseteq M$ be fields. Then $M/K$ is finite if and
only if both $M/L$ and $L/K$ are finite, in which case
\[
  [M:K] = [M:L]\,[L:K].
\]
\end{theorem}

\begin{proof}
Let $\{e_1, \dots, e_m\}$ be a basis of $L$ over $K$ (so
$m = [L:K]$), and let $\{f_1, \dots, f_n\}$ be a basis of $M$ over
$L$ (so $n = [M:L]$). We claim that the set
\[
  \mathcal{B} = \{ e_i f_j : 1 \le i \le m,\; 1 \le j \le n \}
\]
is a basis of $M$ over $K$. This set has $mn$ elements, so the result
will follow.

\textbf{Spanning.} Let $x \in M$. Since $\{f_1, \dots, f_n\}$ is a
basis of $M/L$, we can write
\[
  x = \sum_{j=1}^n \lambda_j f_j, \quad \lambda_j \in L.
\]
Since $\{e_1, \dots, e_m\}$ is a basis of $L/K$, each $\lambda_j$ can
be written as
\[
  \lambda_j = \sum_{i=1}^m a_{ij} e_i, \quad a_{ij} \in K.
\]
Substituting:
\[
  x = \sum_{j=1}^n \Bigl(\sum_{i=1}^m a_{ij} e_i\Bigr) f_j
    = \sum_{i,j} a_{ij}\, e_i f_j.
\]
Hence $\mathcal{B}$ spans $M$ over $K$.

\textbf{Linear independence.} Suppose
\[
  \sum_{i,j} a_{ij}\, e_i f_j = 0, \quad a_{ij} \in K.
\]
Rewriting:
\[
  \sum_{j=1}^n \Bigl(\sum_{i=1}^m a_{ij} e_i\Bigr) f_j = 0.
\]
Set $\mu_j = \sum_{i=1}^m a_{ij} e_i \in L$. Since
$\{f_1, \dots, f_n\}$ is linearly independent over $L$, we have
$\mu_j = 0$ for all $j$. But then
$\sum_{i=1}^m a_{ij} e_i = 0$ with $a_{ij} \in K$, and the
linear independence of $\{e_1, \dots, e_m\}$ over $K$ gives
$a_{ij} = 0$ for all $i, j$. Thus $\mathcal{B}$ is linearly
independent.

For the ``if and only if'' part: if $[M:K]$ is finite, then since
$L$ is a $K$-subspace of $M$ and $M$ is a finite-dimensional
$K$-vector space, $L$ is also finite-dimensional over $K$, so $[L:K]$
is finite. Moreover, any $L$-linearly independent subset of $M$ is also
$K$-linearly independent (since $L \supseteq K$), so $[M:L] \le [M:K]$
and $[M:L]$ is finite.

Conversely, if both $[M:L]$ and $[L:K]$ are finite, then the basis
$\mathcal{B}$ constructed above shows $[M:K] = [M:L][L:K] < \infty$.
\end{proof}

The following diagram illustrates the tower law.
\begin{center}
\begin{tikzcd}[row sep=1.8cm, column sep=0cm]
  M \arrow[d, dash, "{[M:L]}"'] \\
  L \arrow[d, dash, "{[L:K]}"'] \\
  K
\end{tikzcd}
\qquad
\text{gives}\qquad
\begin{tikzcd}[row sep=3.6cm, column sep=0cm]
  M \arrow[d, dash, "{[M:K]=[M:L][L:K]}"'] \\
  K
\end{tikzcd}
\end{center}

\begin{corollary}[Degree divides]
\label{cor:degree-divides}
If $K \subseteq L \subseteq M$ and $[M:K]$ is finite, then $[L:K]$
divides $[M:K]$.
\end{corollary}

\begin{proof}
By the tower law, $[M:K] = [M:L]\cdot[L:K]$.
\end{proof}

\section{Algebraic elements and minimal polynomials}
\label{sec:algebraic-elements}

\begin{definition}[Algebraic and transcendental elements]
\label{def:algebraic-element}
\index{algebraic element}
\index{transcendental element}
Let $L/K$ be a field extension and $\alpha \in L$.
\begin{enumerate}[label=(\roman*)]
  \item $\alpha$ is \textbf{algebraic} over $K$ if there exists a
    nonzero polynomial $f \in K[X]$ such that $f(\alpha) = 0$.
  \item $\alpha$ is \textbf{transcendental} over $K$ if it is not
    algebraic over $K$.
\end{enumerate}
\end{definition}

\begin{definition}[Evaluation homomorphism]
\label{def:eval-hom}
\index{evaluation homomorphism}
For $\alpha \in L$, the \textbf{evaluation homomorphism} is the ring
homomorphism
\[
  \mathrm{ev}_\alpha \colon K[X] \longrightarrow L, \quad
  f(X) \longmapsto f(\alpha).
\]
Its image is the subring $K[\alpha] \subseteq L$, and its kernel
$\Ker(\mathrm{ev}_\alpha)$ is an ideal of $K[X]$.
\end{definition}

Since $K[X]$ is a principal ideal domain\index{principal ideal domain},
$\Ker(\mathrm{ev}_\alpha) = (g)$ for some $g \in K[X]$. If $\alpha$
is algebraic, then $g \ne 0$; if $\alpha$ is transcendental, then
$g = 0$.

\begin{theorem}[Existence and uniqueness of the minimal polynomial]
\label{thm:minimal-poly}
\index{minimal polynomial}
Let $L/K$ be a field extension and let $\alpha \in L$ be algebraic
over $K$. There exists a unique monic polynomial
$m_\alpha = \irr(\alpha, K) \in K[X]$ such that:
\begin{enumerate}[label=(\roman*)]
  \item $m_\alpha(\alpha) = 0$.
  \item If $f \in K[X]$ satisfies $f(\alpha) = 0$, then $m_\alpha
    \mid f$ in $K[X]$.
  \item $m_\alpha$ is irreducible in $K[X]$.
  \item $\Ker(\mathrm{ev}_\alpha) = (m_\alpha)$.
\end{enumerate}
The polynomial $m_\alpha$ is called the \textbf{minimal polynomial}
of $\alpha$ over $K$.
\end{theorem}

\begin{proof}
Since $\alpha$ is algebraic, $I = \Ker(\mathrm{ev}_\alpha)$ is a
nonzero ideal of $K[X]$. Because $K[X]$ is a PID, $I = (m_\alpha)$
for a unique monic generator $m_\alpha \in K[X]$.

\textbf{(i)} Since $\mathrm{ev}_\alpha(m_\alpha) = m_\alpha(\alpha)$
and $m_\alpha \in I = \Ker(\mathrm{ev}_\alpha)$, we have
$m_\alpha(\alpha) = 0$.

\textbf{(ii)} If $f(\alpha) = 0$, then $f \in I = (m_\alpha)$, so
$m_\alpha \mid f$.

\textbf{(iii)} We show $m_\alpha$ is irreducible. Since $m_\alpha$
generates $\Ker(\mathrm{ev}_\alpha)$ and $\alpha$ is algebraic,
$\deg m_\alpha \ge 1$. Suppose $m_\alpha = gh$ with
$g, h \in K[X]$ and $\deg g, \deg h \ge 1$. Then
$0 = m_\alpha(\alpha) = g(\alpha) h(\alpha)$. Since $L$ is a field
(and hence an integral domain), either $g(\alpha) = 0$ or
$h(\alpha) = 0$. Say $g(\alpha) = 0$; then $g \in (m_\alpha)$,
so $m_\alpha \mid g$, which is impossible since
$\deg g < \deg m_\alpha$. This contradiction shows $m_\alpha$ is
irreducible.

\textbf{Uniqueness.} If $m'$ is another monic polynomial satisfying
(i)--(iii), then $m_\alpha \mid m'$ by~(ii) applied to $m_\alpha$,
and $m' \mid m_\alpha$ by the analogous property of $m'$. Since both
are monic, $m_\alpha = m'$.
\end{proof}

\begin{example}[Minimal polynomials]
\label{ex:minimal-polys}
\begin{enumerate}[label=(\roman*)]
  \item $\irr(\sqrt{2}, \Q) = X^2 - 2$. Indeed, $\sqrt{2}$ is a root
    of $X^2 - 2$, and this polynomial is irreducible over $\Q$ by
    Eisenstein's criterion\index{Eisenstein's criterion} with $p = 2$.
  \item $\irr(i, \Q) = X^2 + 1$. The polynomial $X^2+1$ has no
    rational roots and is therefore irreducible over~$\Q$.
  \item $\irr(\sqrt[3]{2}, \Q) = X^3 - 2$, irreducible over $\Q$ by
    Eisenstein with $p = 2$.
  \item $\irr(\sqrt{2}, \Q(\sqrt{2})) = X - \sqrt{2}$, since
    $\sqrt{2} \in \Q(\sqrt{2})$.
\end{enumerate}
\end{example}

\section{The structure of simple extensions}
\label{sec:simple-extensions}

\begin{theorem}[Structure of simple algebraic extensions]
\label{thm:simple-extension-structure}
\index{simple extension}
Let $L/K$ be a field extension and $\alpha \in L$ algebraic over $K$
with minimal polynomial $m_\alpha = \irr(\alpha, K)$ of degree $n$.
Then:
\begin{enumerate}[label=(\roman*)]
  \item There is a $K$-algebra isomorphism
    \[
      K[X]/(m_\alpha) \;\xrightarrow{\;\sim\;}\; K(\alpha),
      \qquad \overline{X} \longmapsto \alpha.
    \]
  \item $K(\alpha) = K[\alpha]$, i.e.\ the field generated by
    $\alpha$ over $K$ equals the ring generated by $\alpha$ over $K$.
  \item $\{1, \alpha, \alpha^2, \dots, \alpha^{n-1}\}$ is a
    $K$-basis of $K(\alpha)$.
  \item $[K(\alpha) : K] = \deg m_\alpha = n$.
\end{enumerate}
\end{theorem}

\begin{proof}
Consider the evaluation homomorphism
$\mathrm{ev}_\alpha \colon K[X] \to L$, $f(X) \mapsto f(\alpha)$.
Its image is $K[\alpha]$ and its kernel is $(m_\alpha)$ by
Theorem~\ref{thm:minimal-poly}. By the first isomorphism theorem
for rings:
\[
  K[X]/(m_\alpha) \;\cong\; K[\alpha].
\]
Since $m_\alpha$ is irreducible in the PID $K[X]$, the ideal
$(m_\alpha)$ is maximal, so $K[X]/(m_\alpha)$ is a field. Therefore
$K[\alpha]$ is a field, and since it is a subfield of $L$ containing
$K$ and $\alpha$, it must be the smallest such field, i.e.\
$K[\alpha] = K(\alpha)$. This proves (i) and (ii).

For~(iii), every element of $K[X]/(m_\alpha)$ is represented by a
polynomial of degree $< n = \deg m_\alpha$ (by the division algorithm),
and these representatives are unique. Under the isomorphism, the
images of $\overline{1}, \overline{X}, \dots, \overline{X^{n-1}}$ are
$1, \alpha, \dots, \alpha^{n-1}$. Since
$\{\overline{1}, \overline{X}, \dots, \overline{X^{n-1}}\}$ is a
$K$-basis of $K[X]/(m_\alpha)$, it follows that
$\{1, \alpha, \dots, \alpha^{n-1}\}$ is a $K$-basis of $K(\alpha)$.
Part~(iv) is immediate from~(iii).
\end{proof}

\begin{example}[$\Q(\sqrt{2})$]
\label{ex:Q-sqrt2}
\index{Q(sqrt2)@$\Q(\sqrt{2})$}
Since $\irr(\sqrt{2}, \Q) = X^2 - 2$ has degree $2$, we have
$[\Q(\sqrt{2}):\Q] = 2$ and
\[
  \Q(\sqrt{2}) = \{ a + b\sqrt{2} : a, b \in \Q \}
  \;\cong\; \Q[X]/(X^2 - 2).
\]
\end{example}

\begin{example}[{$\Q(\sqrt[3]{2})$}]
\label{ex:Q-cbrt2}
\index{Q(cbrt2)@$\Q(\sqrt[3]{2})$}
Since $\irr(\sqrt[3]{2}, \Q) = X^3 - 2$ has degree $3$, we get
$[\Q(\sqrt[3]{2}):\Q] = 3$ and
\[
  \Q(\sqrt[3]{2}) = \{ a + b\sqrt[3]{2} + c\sqrt[3]{4}
    : a, b, c \in \Q \}.
\]
\end{example}

\begin{example}[$\Q(i)$]
\label{ex:Q-i}
\index{Q(i)@$\Q(i)$}
Since $\irr(i, \Q) = X^2 + 1$, we have $[\Q(i):\Q] = 2$ and
\[
  \Q(i) = \{ a + bi : a, b \in \Q \} \cong \Q[X]/(X^2 + 1).
\]
This is the field of \textbf{Gaussian rationals}\index{Gaussian
rationals}.
\end{example}

\section{Finite implies algebraic}
\label{sec:finite-implies-algebraic}

\begin{theorem}[Finite implies algebraic]
\label{thm:finite-implies-algebraic}
\index{finite extension!implies algebraic}
If $L/K$ is a finite extension, then every element of $L$ is algebraic
over $K$.
\end{theorem}

\begin{proof}
Let $n = [L:K]$ and let $\alpha \in L$. Consider the $n+1$ elements
$1, \alpha, \alpha^2, \dots, \alpha^n$ in the $n$-dimensional
$K$-vector space $L$. They must be linearly dependent over $K$, so
there exist $a_0, a_1, \dots, a_n \in K$, not all zero, such that
\[
  a_0 + a_1 \alpha + a_2 \alpha^2 + \cdots + a_n \alpha^n = 0.
\]
The polynomial $f(X) = a_0 + a_1 X + \cdots + a_n X^n \in K[X]$ is
nonzero and satisfies $f(\alpha) = 0$. Hence $\alpha$ is algebraic
over $K$.
\end{proof}

\begin{remark}[Converse is false]
\label{rem:algebraic-not-finite}
The converse is false in general: an extension can be algebraic but
infinite. For example, the algebraic closure $\overline{\Q}/\Q$ is
algebraic but $[\overline{\Q}:\Q] = \infty$.
\end{remark}

\begin{definition}[Algebraic extension]
\label{def:algebraic-extension}
\index{algebraic extension}
A field extension $L/K$ is \textbf{algebraic} if every element of $L$
is algebraic over $K$. Otherwise it is \textbf{transcendental}.
\end{definition}

\begin{theorem}[Algebraic over algebraic is algebraic]
\label{thm:algebraic-tower}
\index{algebraic extension!transitivity}
If $M/L$ is algebraic and $L/K$ is algebraic, then $M/K$ is algebraic.
\end{theorem}

\begin{proof}
Let $\alpha \in M$. We must show that $\alpha$ is algebraic over $K$.
Since $M/L$ is algebraic, $\alpha$ is algebraic over $L$, so there
exists a polynomial
\[
  f(X) = b_0 + b_1 X + \cdots + b_n X^n \in L[X], \quad b_n \ne 0,
\]
with $f(\alpha) = 0$. Consider the intermediate field
$K' = K(b_0, b_1, \dots, b_n)$.

Since $L/K$ is algebraic, each $b_i$ is algebraic over $K$. We
build a tower of finite extensions:
\[
  K \subseteq K(b_0) \subseteq K(b_0, b_1) \subseteq \cdots
  \subseteq K(b_0, \dots, b_n) = K'.
\]
Each step is a simple algebraic extension, hence finite by
Theorem~\ref{thm:simple-extension-structure}(iv). By the tower law
(Theorem~\ref{thm:tower-law}), $[K':K] < \infty$.

Now $\alpha$ is algebraic over $K' \supseteq L$ (since
$f \in K'[X]$ and $f(\alpha) = 0$), so
$[K'(\alpha):K'] \le n < \infty$. Applying the tower law once more:
\[
  [K'(\alpha) : K] = [K'(\alpha):K']\cdot[K':K] < \infty.
\]
Since $K \subseteq K'(\alpha)$ is a finite extension, $\alpha$ is
algebraic over $K$ by Theorem~\ref{thm:finite-implies-algebraic}.
\end{proof}

\begin{corollary}[Set of algebraic elements is a field]
\label{cor:algebraic-elements-field}
Let $L/K$ be a field extension. The set
\[
  \overline{K}_L = \{ \alpha \in L : \alpha \text{ is algebraic over }
  K \}
\]
is a subfield of $L$ containing $K$.
\end{corollary}

\begin{proof}
Clearly $K \subseteq \overline{K}_L$. We must show that if
$\alpha, \beta \in \overline{K}_L$ with $\beta \ne 0$, then
$\alpha + \beta$, $\alpha - \beta$, $\alpha\beta$, and
$\alpha/\beta$ all lie in $\overline{K}_L$. These elements all
belong to $K(\alpha, \beta)$, and
\[
  [K(\alpha,\beta):K] \le [K(\alpha):K]\cdot[K(\alpha,\beta):K(\alpha)]
  < \infty,
\]
since $\alpha$ is algebraic over $K$ and $\beta$ is algebraic over
$K(\alpha)$ (because it is algebraic over $K \subseteq K(\alpha)$).
By Theorem~\ref{thm:finite-implies-algebraic}, every element of
$K(\alpha,\beta)$ is algebraic over $K$.
\end{proof}

\section{Transcendental extensions}
\label{sec:transcendental-extensions}

\begin{definition}[Purely transcendental extension]
\label{def:purely-transcendental}
\index{transcendental extension}
\index{purely transcendental extension}
Let $L/K$ be a field extension and $\alpha \in L$ transcendental
over $K$. The evaluation homomorphism
$\mathrm{ev}_\alpha \colon K[X] \to L$ is injective (since
$\Ker = (0)$), so it extends to an injective field homomorphism
$K(X) \hookrightarrow L$, giving an isomorphism
$K(\alpha) \cong K(X)$, the field of rational functions.
\end{definition}

\begin{remark}[Transcendence degree]
\label{rem:transcendence-degree}
\index{transcendence degree}
The study of transcendental extensions leads to the notion of
\emph{transcendence degree}, which measures how many ``independent
transcendentals'' are needed. The transcendence degree of
$\R/\Q$ is $|\R|$ (uncountable). We will not develop this theory
in detail, as our focus is on algebraic extensions and Galois theory.
\end{remark}

\section{Characteristic and the Frobenius endomorphism}
\label{sec:characteristic-frobenius}

\begin{definition}[Characteristic]
\label{def:characteristic}
\index{characteristic!of a field}
The \textbf{characteristic} of a field $K$, denoted $\chr(K)$, is
the smallest positive integer $p$ such that
$\underbrace{1 + 1 + \cdots + 1}_{p} = 0$ in $K$. If no such integer
exists, we set $\chr(K) = 0$.
\end{definition}

\begin{proposition}[Characteristic is $0$ or prime]
\label{prop:char-prime}
\index{characteristic!is prime}
The characteristic of a field is either $0$ or a prime number.
\end{proposition}

\begin{proof}
The unique ring homomorphism $\Z \to K$, $n \mapsto n \cdot 1_K$, has
kernel $(p)$ for some $p \ge 0$ (since $\Z$ is a PID). If $p \ne 0$,
the image $\Z/(p)$ embeds in $K$, so $\Z/(p)$ is an integral domain,
which forces $p$ to be prime.
\end{proof}

\begin{definition}[Prime subfield]
\label{def:prime-subfield}
\index{prime subfield}
The \textbf{prime subfield} of $K$ is the intersection of all
subfields of $K$. It is isomorphic to $\Q$ if $\chr(K) = 0$, and
to $\F_p$ if $\chr(K) = p$.
\end{definition}

\begin{theorem}[Frobenius endomorphism]
\label{thm:frobenius}
\index{Frobenius endomorphism}
Let $K$ be a field of characteristic $p > 0$. The map
\[
  \varphi \colon K \longrightarrow K, \quad x \longmapsto x^p,
\]
is an injective field homomorphism, called the \textbf{Frobenius
endomorphism}.
\end{theorem}

\begin{proof}
We verify the homomorphism properties.

\textbf{Multiplicativity.} For all $x, y \in K$:
$\varphi(xy) = (xy)^p = x^p y^p = \varphi(x)\varphi(y)$.

\textbf{Additivity.} For all $x, y \in K$:
$\varphi(x + y) = (x+y)^p$. By the binomial theorem,
\[
  (x+y)^p = \sum_{k=0}^{p} \binom{p}{k} x^k y^{p-k}.
\]
For $1 \le k \le p-1$, the binomial coefficient
$\binom{p}{k} = \frac{p!}{k!(p-k)!}$ is divisible by $p$ (since $p$
is prime and does not divide $k!$ or $(p-k)!$ for $0 < k < p$).
Hence $\binom{p}{k} \equiv 0$ in $K$ for $1 \le k \le p-1$, and
\[
  (x+y)^p = x^p + y^p = \varphi(x) + \varphi(y).
\]

\textbf{$\varphi(1) = 1$.} This is clear since $1^p = 1$.

\textbf{Injectivity.} Since $K$ is a field and $\varphi$ is a nonzero
ring homomorphism, $\Ker(\varphi)$ is an ideal of $K$. The only
ideals of a field are $(0)$ and $K$ itself. Since
$\varphi(1) = 1 \ne 0$, we have $\Ker(\varphi) = (0)$, so $\varphi$
is injective.
\end{proof}

\begin{remark}[Frobenius and finite fields]
\label{rem:frobenius-finite}
If $K$ is a finite field, then $\varphi$ is also surjective (an
injective map from a finite set to itself is a bijection), so
$\varphi$ is an automorphism. For infinite fields of characteristic
$p$, the Frobenius need not be surjective. For instance, in
$\F_p(t)$ (rational functions over $\F_p$), the element $t$ is not
a $p$-th power.
\end{remark}

\section{Finite fields: first examples}
\label{sec:finite-fields-intro}

\begin{example}[Finite fields of small order]
\label{ex:finite-fields-small}
\index{finite field!small examples}
\begin{enumerate}[label=(\roman*)]
  \item $\F_2 = \{0, 1\}$, the field with two elements.
  \item $\F_4$: since $4 = 2^2$, we seek an irreducible polynomial of
    degree $2$ over $\F_2$. The polynomial $X^2 + X + 1$ is
    irreducible over $\F_2$ (it has no roots in $\F_2$), so
    \[
      \F_4 \cong \F_2[X]/(X^2 + X + 1) = \{0, 1, \alpha, 1+\alpha\},
    \]
    where $\alpha^2 = \alpha + 1$.
  \item $\F_8 \cong \F_2[X]/(X^3 + X + 1)$, since $X^3 + X + 1$ is
    irreducible over $\F_2$.
\end{enumerate}
\end{example}

The following extension tower illustrates the relationship between
these fields.

\begin{center}
\begin{tikzcd}[row sep=1.5cm, column sep=0cm]
  \F_8 \\
  & \F_4 \arrow[ul, dash, "3"'] \\
  & \F_2 \arrow[u, dash, "2"'] \arrow[uul, dash, "3"]
\end{tikzcd}
\end{center}

Note that $\F_4$ is \emph{not} a subfield of $\F_8$ (since $2 \nmid 3$),
which is why the diagram shows no direct inclusion arrow from $\F_4$ to
$\F_8$.

\section{Exercises}
\label{sec:ch8-exercises}

\begin{exercise}[Degree computations]
\label{exer:degree-computations}
Determine $[\Q(\sqrt{3}, \sqrt{5}) : \Q]$ and find a basis for
$\Q(\sqrt{3}, \sqrt{5})$ over $\Q$.
\textit{Hint:} Use the tower law and show that $\sqrt{5} \notin
\Q(\sqrt{3})$.
\end{exercise}

\begin{exercise}[Minimal polynomial over an extension]
\label{exer:min-poly-ext}
Find $\irr(\sqrt{2}+\sqrt{3}, \Q)$ and show that it has degree $4$.
Deduce that $\Q(\sqrt{2}+\sqrt{3}) = \Q(\sqrt{2}, \sqrt{3})$.
\end{exercise}

\begin{exercise}[Algebraic element in a tower]
\label{exer:algebraic-tower}
Let $K \subseteq L \subseteq M$ be field extensions. Show that if
$\alpha \in M$ is algebraic over $K$, then it is algebraic over $L$,
and $\deg(\irr(\alpha, L)) \le \deg(\irr(\alpha, K))$.
\end{exercise}

\begin{exercise}[Degree and divisibility]
\label{exer:degree-divides}
Let $\alpha$ be algebraic over $K$ with $\deg(\irr(\alpha,K)) = n$.
Show that if $f \in K[X]$ satisfies $f(\alpha) = 0$ and
$\deg f = n$, then $f$ is a scalar multiple of $\irr(\alpha,K)$.
\end{exercise}

\begin{exercise}[Simple extensions of finite fields]
\label{exer:simple-ext-finite}
Show that $\F_{p^n}$ can be obtained as a simple extension
$\F_p(\alpha)$ for some $\alpha$.
\textit{Hint:} Use the fact that $\F_{p^n}^{\times}$ is cyclic
(proved later in Theorem~\ref{thm:finite-field-mult-cyclic}).
\end{exercise}

\begin{exercise}[Compositum]
\label{exer:compositum}
\index{compositum}
Let $L_1/K$ and $L_2/K$ be finite extensions inside a common
overfield $M$. Define the \textbf{compositum} $L_1 L_2$ as the
smallest subfield of $M$ containing both $L_1$ and $L_2$. Show that
\[
  [L_1 L_2 : K] \le [L_1 : K]\,[L_2 : K],
\]
with equality when $\gcd([L_1:K], [L_2:K]) = 1$.
\end{exercise}

\begin{exercise}[Transcendence of $e$ and $\pi$]
\label{exer:transcendence}
State (without proof) the Lindemann--Weierstrass theorem and use it
to show that $e$ and $\pi$ are transcendental over $\Q$.
\end{exercise}

\begin{exercise}[Frobenius endomorphism]
\label{exer:frobenius}
Let $K$ be a field of characteristic $p > 0$. For $n \ge 1$, define
$\varphi^n \colon K \to K$ by $\varphi^n(x) = x^{p^n}$. Show that
$\varphi^n$ is a field homomorphism.
\end{exercise}

\begin{exercise}[Subfield criterion]
\label{exer:subfield-criterion}
Let $L/K$ be a finite extension with $[L:K] = p$, a prime. Show that
there are no intermediate fields strictly between $K$ and $L$.
\end{exercise}

\begin{exercise}[Algebraic numbers form a field]
\label{exer:alg-numbers-field}
Let $\overline{\Q}$ denote the set of all complex numbers that are
algebraic over $\Q$. Verify directly (without using
Corollary~\ref{cor:algebraic-elements-field}) that $\overline{\Q}$
is a subfield of $\C$.
\end{exercise}

\begin{exercise}[Extension of degree $2$]
\label{exer:deg-2-extension}
Let $K$ be a field with $\chr(K) \ne 2$ and $L/K$ an extension with
$[L:K]=2$. Show that $L = K(\sqrt{d})$ for some $d \in K \setminus
K^2$.
\end{exercise}

\begin{exercise}[Infinite algebraic extension]
\label{exer:infinite-algebraic}
Consider $L = \Q(\sqrt{2}, \sqrt{3}, \sqrt{5}, \sqrt{7}, \dots)$
(adjoining $\sqrt{p}$ for every prime $p$). Show that $L/\Q$ is
algebraic but not finite.
\end{exercise}

\subsection*{Chapter summary}
\addcontentsline{toc}{section}{Chapter summary}

\begin{itemize}[leftmargin=2em]
  \item A \textbf{field extension} $L/K$ makes $L$ a $K$-vector space;
    the degree $[L:K]$ is its dimension.
  \item The \textbf{tower law} gives
    $[M:K] = [M:L][L:K]$.
  \item An algebraic element $\alpha$ has a unique monic irreducible
    \textbf{minimal polynomial} $\irr(\alpha,K)$; the simple extension
    $K(\alpha) \cong K[X]/(\irr(\alpha,K))$ has degree equal to the
    degree of $\irr(\alpha,K)$.
  \item Finite $\Rightarrow$ algebraic, and algebraic over algebraic
    is algebraic.
  \item The \textbf{Frobenius endomorphism} $x \mapsto x^p$ is an
    injective field homomorphism in characteristic $p$.
\end{itemize}


%% ══════════════════════════════════════════════════════════
%% CHAPTER 9: SPLITTING FIELDS
%% ══════════════════════════════════════════════════════════
\chapter{Splitting Fields}
\label{ch:splitting-fields}
\index{splitting field}

\section{Definition and first examples}
\label{sec:splitting-def}

\begin{definition}[Root of a polynomial in an extension]
\label{def:root-extension}
\index{root!in an extension}
Let $K$ be a field and $f \in K[X]$ a polynomial of degree $n \ge 1$.
An element $\alpha$ in some extension $L/K$ is a \textbf{root} of $f$
if $f(\alpha) = 0$.
\end{definition}

\begin{definition}[Polynomial splits]
\label{def:poly-splits}
\index{splitting!of a polynomial}
We say that $f \in K[X]$ \textbf{splits} (or \textbf{splits
completely}) over an extension field $L \supseteq K$ if $f$ factors
as a product of linear factors in $L[X]$:
\[
  f(X) = a \prod_{i=1}^{n} (X - \alpha_i), \quad
  \alpha_i \in L,\; a \in K^{\times}.
\]
\end{definition}

\begin{definition}[Splitting field]
\label{def:splitting-field}
\index{splitting field!definition}
Let $f \in K[X]$ with $\deg f \ge 1$. A field $L \supseteq K$ is a
\textbf{splitting field} of $f$ over $K$ if:
\begin{enumerate}[label=(\roman*)]
  \item $f$ splits completely over $L$: writing
    $f(X) = a\prod_{i=1}^n (X - \alpha_i)$ with $\alpha_i \in L$;
  \item $L$ is generated over $K$ by the roots of $f$:
    $L = K(\alpha_1, \dots, \alpha_n)$.
\end{enumerate}
Condition~(ii) ensures that $L$ is the \emph{smallest} extension of $K$
over which $f$ splits.
\end{definition}

\begin{remark}[One root versus all roots]
\label{rem:one-vs-all}
Given an irreducible polynomial $p(X) \in K[X]$ of degree $n$, the
quotient $K[X]/(p(X))$ is a field extension of $K$ of degree $n$ in
which $p$ has \emph{at least one} root (namely $\overline{X}$), but
$p$ need not split completely there. For example, $X^3 - 2 \in
\Q[X]$ has one real root $\sqrt[3]{2}$, but $\Q(\sqrt[3]{2})
\subseteq \R$ contains no complex cube roots of $2$, so $X^3 - 2$
does not split over $\Q(\sqrt[3]{2})$.
\end{remark}

\begin{example}[Splitting fields]
\label{ex:splitting-fields}
\index{splitting field!examples}
\begin{enumerate}[label=(\roman*)]
  \item \textbf{$X^2 + 1$ over $\R$:} The roots are $\pm i$, so the
    splitting field is $\R(i) = \C$, with $[\C:\R] = 2$.
  \item \textbf{$X^2 - 2$ over $\Q$:} The roots are $\pm\sqrt{2}$,
    and $\Q(\sqrt{2}, -\sqrt{2}) = \Q(\sqrt{2})$, so the splitting
    field is $\Q(\sqrt{2})$ with degree $2$ over $\Q$.
  \item \textbf{$X^3 - 2$ over $\Q$:}\label{ex:splitting-X3-2}
    The three roots are $\sqrt[3]{2}$, $\sqrt[3]{2}\,\omega$,
    $\sqrt[3]{2}\,\omega^2$, where $\omega = e^{2\pi i/3}$. The
    splitting field is
    $\Q(\sqrt[3]{2}, \omega) \subseteq \C$.

    We compute the degree using the tower
    \begin{center}
    \begin{tikzcd}[row sep=1.5cm, column sep=0cm]
      \Q(\sqrt[3]{2}, \omega)
        \arrow[d, dash, "2"'] \\
      \Q(\sqrt[3]{2})
        \arrow[d, dash, "3"'] \\
      \Q
    \end{tikzcd}
    \end{center}
    Indeed, $[\Q(\sqrt[3]{2}):\Q] = 3$, and
    $\omega \notin \Q(\sqrt[3]{2}) \subseteq \R$, so $X^2 + X + 1$
    remains irreducible over $\Q(\sqrt[3]{2})$ (it has degree $2$
    and no real roots). Hence
    $[\Q(\sqrt[3]{2},\omega):\Q(\sqrt[3]{2})] = 2$ and
    \[
      [\Q(\sqrt[3]{2}, \omega) : \Q] = 3 \cdot 2 = 6.
    \]
\end{enumerate}
\end{example}

\begin{example}[Cyclotomic splitting fields]
\label{ex:cyclotomic-splitting}
\index{cyclotomic polynomial}
The $n$-th \textbf{cyclotomic polynomial} $\Phi_n(X) \in \Z[X]$ is
the minimal polynomial of a primitive $n$-th root of unity over $\Q$.
Its splitting field over $\Q$ is the \textbf{cyclotomic field}
$\Q(\zeta_n)$, where $\zeta_n = e^{2\pi i/n}$, and
$[\Q(\zeta_n):\Q] = \varphi(n)$ (Euler's totient function).
\end{example}

\section{Existence of splitting fields}
\label{sec:splitting-existence}

\begin{theorem}[Existence of splitting fields]
\label{thm:splitting-existence}
\index{splitting field!existence}
For every field $K$ and every polynomial $f \in K[X]$ with
$\deg f = n \ge 1$, there exists a splitting field $L$ of $f$ over
$K$, with $[L:K] \le n!$.
\end{theorem}

\begin{proof}
We proceed by induction on $n = \deg f$.

\textbf{Base case ($n = 1$).} If $\deg f = 1$, then $f$ is already
linear, $f(X) = a(X - \alpha)$ with $\alpha \in K$, and $L = K$ is a
splitting field with $[K:K] = 1 \le 1!$.

\textbf{Inductive step.} Assume the result holds for all fields and
all polynomials of degree $< n$. Let $f \in K[X]$ have degree $n$.
Let $p(X) \in K[X]$ be an irreducible factor of $f$ with
$d = \deg p \ge 1$. Set
\[
  K_1 = K[X]/(p(X)).
\]
Then $K_1/K$ is a field extension with $[K_1:K] = d \le n$, and $p$
has a root $\alpha = \overline{X}$ in $K_1$. In $K_1[X]$, we can
write $f(X) = (X - \alpha)\,g(X)$ for some $g \in K_1[X]$ with
$\deg g = n - 1$.

By the inductive hypothesis, there exists a splitting field $L$ of
$g$ over $K_1$ with $[L:K_1] \le (n-1)!$. In $L[X]$, the polynomial
$g$ splits completely, and therefore so does $f(X) = (X-\alpha)g(X)$.

Let $\alpha_1 = \alpha, \alpha_2, \dots, \alpha_n$ be the roots of
$f$ in $L$. The subfield $K(\alpha_1, \dots, \alpha_n) \subseteq L$
is a splitting field of $f$ over $K$. (In fact $L$ itself equals
$K(\alpha_1, \dots, \alpha_n)$ by minimality of $L$ as a splitting
field of $g$ over $K_1 = K(\alpha_1)$.)

For the degree bound:
\[
  [L:K] = [L:K_1]\cdot[K_1:K] \le (n-1)!\cdot d \le (n-1)!\cdot n
  = n!. \qedhere
\]
\end{proof}

\section{Extension of isomorphisms}
\label{sec:extension-iso}

The key technical tool for proving uniqueness of splitting fields is
the following lemma about extending field isomorphisms.

\begin{lemma}[Extension of isomorphisms]
\label{lem:extension-iso}
\index{extension of isomorphisms}
Let $\sigma \colon K \xrightarrow{\sim} K'$ be a field isomorphism.
Let $f \in K[X]$ be irreducible and let $f' = \sigma(f) \in K'[X]$
(obtained by applying $\sigma$ to each coefficient). Suppose
$\alpha$ is a root of $f$ in some extension $L/K$ and $\beta$ is a
root of $f'$ in some extension $L'/K'$. Then $\sigma$ extends to an
isomorphism
\[
  \tilde{\sigma} \colon K(\alpha) \xrightarrow{\;\sim\;} K'(\beta)
\]
with $\tilde{\sigma}(\alpha) = \beta$ and
$\tilde{\sigma}|_K = \sigma$.
\end{lemma}

\begin{proof}
Since $f$ is irreducible over $K$, it is (up to a unit) the minimal
polynomial of $\alpha$ over $K$. By
Theorem~\ref{thm:simple-extension-structure}, we have isomorphisms
\[
  K(\alpha) \cong K[X]/(f), \qquad K'(\beta) \cong K'[X]/(f').
\]
More precisely, the evaluation homomorphisms give
\[
  \psi \colon K[X]/(f) \xrightarrow{\sim} K(\alpha), \quad
  \overline{X} \mapsto \alpha,
\]
\[
  \psi' \colon K'[X]/(f') \xrightarrow{\sim} K'(\beta), \quad
  \overline{X} \mapsto \beta.
\]
The isomorphism $\sigma \colon K \to K'$ induces an isomorphism
$\hat{\sigma} \colon K[X] \to K'[X]$ by applying $\sigma$ to each
coefficient and fixing $X$. Since $\hat{\sigma}(f) = f'$,
$\hat{\sigma}$ maps the ideal $(f)$ onto $(f')$ and therefore induces
an isomorphism
\[
  \bar{\sigma} \colon K[X]/(f) \xrightarrow{\;\sim\;} K'[X]/(f').
\]
The desired extension is then
\[
  \tilde{\sigma} = \psi' \circ \bar{\sigma} \circ \psi^{-1}
  \colon K(\alpha) \xrightarrow{\;\sim\;} K'(\beta).
\]
By construction, $\tilde{\sigma}(\alpha) = \psi'(\bar{\sigma}(
\overline{X})) = \psi'(\overline{X}) = \beta$, and for $a \in K$,
$\tilde{\sigma}(a) = \psi'(\bar{\sigma}(\overline{a})) =
\psi'(\overline{\sigma(a)}) = \sigma(a)$.
\end{proof}

\section{Uniqueness of splitting fields}
\label{sec:splitting-uniqueness}

\begin{theorem}[Uniqueness of splitting fields]
\label{thm:splitting-uniqueness}
\index{splitting field!uniqueness}
Let $\sigma \colon K \xrightarrow{\sim} K'$ be a field isomorphism
and let $f \in K[X]$ with $f' = \sigma(f) \in K'[X]$. Let $L$ be a
splitting field of $f$ over $K$ and $L'$ a splitting field of $f'$
over $K'$. Then $\sigma$ extends to an isomorphism
$\tilde{\sigma} \colon L \xrightarrow{\sim} L'$.

In particular (taking $K = K'$, $\sigma = \id_K$), any two splitting
fields of $f$ over $K$ are $K$-isomorphic.
\end{theorem}

\begin{proof}
We proceed by induction on $n = [L:K]$.

\textbf{Base case ($n = 1$).} If $[L:K] = 1$, then $L = K$ and $f$
already splits over $K$. Then $f' = \sigma(f)$ already splits over
$K'$, so $L' = K'$ (by minimality). The map $\sigma$ itself is the
desired isomorphism.

\textbf{Inductive step.} Assume $n > 1$ and that the result holds for
all splitting fields of smaller degree over their base.

Since $[L:K] > 1$, the polynomial $f$ has an irreducible factor
$p \in K[X]$ with $\deg p \ge 2$. Let $\alpha \in L$ be a root of $p$
and set $p' = \sigma(p) \in K'[X]$. Since $f'$ splits over $L'$ and
$p' \mid f'$, the polynomial $p'$ has a root $\beta \in L'$.

By Lemma~\ref{lem:extension-iso}, $\sigma$ extends to an isomorphism
\[
  \sigma_1 \colon K(\alpha) \xrightarrow{\;\sim\;} K'(\beta),
  \quad \sigma_1(\alpha) = \beta.
\]

Now $L$ is a splitting field of $f$ over $K(\alpha)$ (since
$L = K(\alpha_1, \dots, \alpha_n) = K(\alpha)(\alpha_2, \dots,
\alpha_n)$ and $f$ splits over $L$), and similarly $L'$ is a
splitting field of $f'$ over $K'(\beta)$. Moreover,
\[
  [L : K(\alpha)] = \frac{[L:K]}{[K(\alpha):K]}
  = \frac{n}{[K(\alpha):K]} < n,
\]
since $[K(\alpha):K] = \deg p \ge 2$.

By the inductive hypothesis applied to $\sigma_1 \colon K(\alpha) \to
K'(\beta)$, there exists an isomorphism
$\tilde{\sigma} \colon L \xrightarrow{\sim} L'$ extending $\sigma_1$.
Since $\sigma_1$ extends $\sigma$, so does $\tilde{\sigma}$.
\end{proof}

\begin{corollary}[Splitting fields are unique up to isomorphism]
\label{cor:splitting-unique-isom}
Any two splitting fields of $f \in K[X]$ over $K$ are $K$-isomorphic.
The isomorphism need not be unique.
\end{corollary}

\section{Multiple roots and separability}
\label{sec:separability}

\begin{definition}[Multiple root]
\label{def:multiple-root}
\index{multiple root}
Let $f \in K[X]$ and let $L$ be a splitting field of $f$ over $K$.
A root $\alpha \in L$ of $f$ is a \textbf{multiple root} (or
\textbf{repeated root}) if $(X - \alpha)^2 \mid f$ in $L[X]$. If
$(X - \alpha)^m \mid f$ but $(X - \alpha)^{m+1} \nmid f$, we say
$\alpha$ has \textbf{multiplicity}\index{multiplicity} $m$.
A root of multiplicity $1$ is called \textbf{simple}.
\end{definition}

\begin{definition}[Formal derivative]
\label{def:formal-derivative}
\index{formal derivative}
For $f = \sum_{k=0}^n a_k X^k \in K[X]$, the \textbf{formal
derivative} is
\[
  f' = \sum_{k=1}^n k\, a_k X^{k-1} \in K[X].
\]
The usual rules hold: $(f+g)' = f' + g'$, $(fg)' = f'g + fg'$,
$(cf)' = cf'$ for $c \in K$.
\end{definition}

\begin{proposition}[Multiple roots and the derivative]
\label{prop:multiple-root-derivative}
\index{multiple root!and derivative}
Let $f \in K[X]$ be nonconstant and let $\alpha$ be a root of $f$ in
some extension of $K$. Then $\alpha$ is a multiple root of $f$ if and
only if $f'(\alpha) = 0$.
\end{proposition}

\begin{proof}
Write $f(X) = (X - \alpha)^m g(X)$ in the splitting field, where
$g(\alpha) \ne 0$ and $m \ge 1$.

$(\Rightarrow)$ If $m \ge 2$, then
$f'(X) = m(X-\alpha)^{m-1}g(X) + (X-\alpha)^m g'(X)$, so
$f'(\alpha) = 0$.

$(\Leftarrow)$ If $m = 1$, then
$f(X) = (X-\alpha)g(X)$ and
$f'(X) = g(X) + (X-\alpha)g'(X)$, giving
$f'(\alpha) = g(\alpha) \ne 0$.
\end{proof}

\begin{corollary}[Criterion for no multiple roots]
\label{cor:no-multiple-roots}
A polynomial $f \in K[X]$ has no multiple roots (in any extension) if
and only if $\gcd(f, f') = 1$ in $K[X]$.
\end{corollary}

\begin{proof}
The polynomial $f$ has a multiple root $\alpha$ (in some extension) if
and only if $f$ and $f'$ share the root $\alpha$, which happens if and
only if some irreducible factor of $f$ divides both $f$ and $f'$,
i.e.\ $\gcd(f, f') \ne 1$.
\end{proof}

\begin{definition}[Separable polynomial]
\label{def:separable-poly}
\index{separable polynomial}
A polynomial $f \in K[X]$ is \textbf{separable} if it has no multiple
roots in any extension of $K$, equivalently if $\gcd(f, f') = 1$.
An irreducible polynomial $f$ is separable if and only if $f' \ne 0$.
\end{definition}

\begin{definition}[Separable extension]
\label{def:separable-extension}
\index{separable extension}
An algebraic extension $L/K$ is \textbf{separable} if the minimal
polynomial $\irr(\alpha, K)$ is separable for every $\alpha \in L$.
\end{definition}

\begin{theorem}[Characteristic zero implies separable]
\label{thm:char0-separable}
\index{separable!in characteristic zero}
If $\chr(K) = 0$, then every irreducible polynomial in $K[X]$ is
separable. Hence every algebraic extension of $K$ is separable.
\end{theorem}

\begin{proof}
Let $f \in K[X]$ be irreducible with $\deg f = n \ge 1$. Write
$f = a_n X^n + a_{n-1}X^{n-1} + \cdots + a_0$ with $a_n \ne 0$. The
formal derivative is
\[
  f' = n a_n X^{n-1} + (n-1)a_{n-1}X^{n-2} + \cdots + a_1.
\]
Since $\chr(K) = 0$, the integer $n$ is nonzero in $K$, and
$a_n \ne 0$, so $n a_n \ne 0$ in $K$. Therefore $\deg f' = n - 1$,
in particular $f' \ne 0$.

Since $f$ is irreducible and $\deg f' < \deg f$, the polynomial $f$
cannot divide $f'$. But $f$ is irreducible, so $\gcd(f, f') = 1$
(the gcd must be $1$ or an associate of $f$; it cannot be an
associate of $f$ since $f \nmid f'$). By
Corollary~\ref{cor:no-multiple-roots}, $f$ is separable.
\end{proof}

\begin{proposition}[Inseparable irreducible polynomials in
  characteristic $p$]
\label{prop:inseparable-char-p}
\index{inseparable polynomial}
Let $K$ be a field with $\chr(K) = p > 0$ and let $f \in K[X]$ be
irreducible. Then $f$ is inseparable if and only if $f' = 0$, which
happens if and only if $f(X) = g(X^p)$ for some $g \in K[X]$.
\end{proposition}

\begin{proof}
If $f$ is inseparable, then $\gcd(f, f') \ne 1$. Since $f$ is
irreducible and $\deg f' < \deg f$, we cannot have $f \mid f'$
unless $f' = 0$. Since the gcd of $f$ and $f'$ divides $f$ and $f$
is irreducible, $\gcd(f,f')$ is either $1$ or an associate of $f$.
If $\gcd(f,f') \ne 1$, then $f \mid f'$, which forces $f' = 0$.

Conversely, if $f' = 0$, then $\gcd(f,f') = f \ne 1$, so $f$ is
inseparable.

Now $f' = 0$ means that $k a_k = 0$ for all $k \ge 1$. In
characteristic $p$, this holds if and only if $a_k = 0$ whenever
$p \nmid k$, i.e.\ $f$ involves only powers of $X$ that are
multiples of $p$: $f(X) = \sum_j b_j X^{pj} = g(X^p)$ where
$g(Y) = \sum_j b_j Y^j$.
\end{proof}

\begin{definition}[Perfect field]
\label{def:perfect-field}
\index{perfect field}
A field $K$ is \textbf{perfect} if every irreducible polynomial in
$K[X]$ is separable. Equivalently:
\begin{itemize}
  \item If $\chr(K) = 0$, then $K$ is automatically perfect.
  \item If $\chr(K) = p > 0$, then $K$ is perfect if and only if
    the Frobenius $\varphi \colon K \to K$, $x \mapsto x^p$, is
    surjective (i.e.\ every element of $K$ has a $p$-th root in $K$).
\end{itemize}
\end{definition}

\begin{proposition}[Finite fields are perfect]
\label{prop:finite-perfect}
Every finite field is perfect.
\end{proposition}

\begin{proof}
Let $K$ be a finite field of characteristic $p$. The Frobenius
$\varphi \colon K \to K$ is injective
(Theorem~\ref{thm:frobenius}). Since $K$ is finite, an injective
map from $K$ to itself is surjective. Hence $\varphi$ is surjective,
and $K$ is perfect.
\end{proof}

\section{Exercises}
\label{sec:ch9-exercises}

\begin{exercise}[Splitting field of $X^4 - 2$]
\label{exer:split-X4-2}
\index{splitting field!of $X^4-2$}
Find the splitting field of $X^4 - 2$ over $\Q$ and compute its
degree over $\Q$.
\end{exercise}

\begin{exercise}[Splitting field over a finite field]
\label{exer:split-finite-field}
Find the splitting field of $X^4 + X + 1$ over $\F_2$ and its degree.
\end{exercise}

\begin{exercise}[Adjoining one root versus all roots]
\label{exer:one-root-all}
Give an example of an irreducible polynomial $f \in \Q[X]$ of degree
$3$ such that adjoining one root of $f$ does not produce the splitting
field. Compute $[K(\alpha):\Q]$ and $[\text{splitting field}:\Q]$.
\end{exercise}

\begin{exercise}[Splitting field of a product]
\label{exer:split-product}
Let $f, g \in K[X]$. Show that the splitting field of $fg$ over $K$
is the compositum of the splitting fields of $f$ and $g$.
\end{exercise}

\begin{exercise}[Counting embeddings]
\label{exer:counting-embeddings}
Let $f \in K[X]$ be irreducible of degree $n$ and let $L$ be a
splitting field of $f$. Show that the number of $K$-embeddings
$K[X]/(f) \hookrightarrow L$ equals the number of distinct roots of
$f$ in $L$, which is $n$ if $f$ is separable.
\end{exercise}

\begin{exercise}[Separable implies distinct roots]
\label{exer:separable-distinct}
Let $f \in K[X]$ be a separable polynomial of degree $n$. Prove that
$f$ has exactly $n$ distinct roots in any splitting field.
\end{exercise}

\begin{exercise}[Inseparable example]
\label{exer:inseparable-example}
Let $K = \F_p(t)$ (the field of rational functions over $\F_p$).
Show that $X^p - t \in K[X]$ is irreducible but inseparable.
\end{exercise}

\begin{exercise}[Derivative and gcd]
\label{exer:derivative-gcd}
Let $f \in K[X]$ be a nonconstant polynomial. Show that
$f / \gcd(f, f')$ is a separable polynomial with the same roots as
$f$ (but all with multiplicity one).
\end{exercise}

\begin{exercise}[Splitting field degree bound]
\label{exer:split-degree-bound}
Show that if $f \in K[X]$ is irreducible of degree $n$, then the
degree of the splitting field $L$ of $f$ over $K$ satisfies
$n \mid [L:K]$ and $[L:K] \mid n!$.
\end{exercise}

\begin{exercise}[Cyclotomic polynomials]
\label{exer:cyclotomic}
\index{cyclotomic polynomial!exercises}
\begin{enumerate}[label=(\alph*)]
  \item Compute $\Phi_n(X)$ for $n = 1, 2, 3, 4, 5, 6, 8, 12$.
  \item Show that $\Phi_p(X) = X^{p-1} + X^{p-2} + \cdots + X + 1$
    for $p$ prime.
  \item Prove that $\Phi_p(X)$ is irreducible over $\Q$ for $p$ prime.
    \textit{Hint:} Substitute $X \mapsto X+1$ and apply Eisenstein.
\end{enumerate}
\end{exercise}

\begin{exercise}[Separability in characteristic $p$]
\label{exer:sep-char-p}
Let $K$ be a field of characteristic $p > 0$ and $f \in K[X]$
irreducible. Show that there exists a unique integer $e \ge 0$ and a
unique separable irreducible polynomial $g \in K[X]$ such that
$f(X) = g(X^{p^e})$.
\end{exercise}

\begin{exercise}[Splitting field is always finite]
\label{exer:split-always-finite}
Let $f \in K[X]$ with $\deg f = n \ge 1$. Prove that the splitting
field of $f$ over $K$ is always a finite extension of $K$, with degree
at most $n!$.
\end{exercise}

\subsection*{Chapter summary}
\addcontentsline{toc}{section}{Chapter summary}

\begin{itemize}[leftmargin=2em]
  \item A \textbf{splitting field} of $f \in K[X]$ is the smallest
    extension of $K$ over which $f$ factors into linear factors.
  \item Splitting fields \textbf{exist} (by adjoining roots
    iteratively) and are \textbf{unique up to $K$-isomorphism}
    (by the extension-of-isomorphisms lemma).
  \item The degree of the splitting field of $f$ divides $(\deg f)!$.
  \item A polynomial is \textbf{separable} if and only if
    $\gcd(f, f') = 1$, equivalently it has no multiple roots.
  \item In characteristic $0$, every irreducible polynomial is
    separable. In characteristic $p$, an irreducible $f$ is
    inseparable iff $f(X) = g(X^p)$.
  \item \textbf{Perfect fields} are those where every irreducible is
    separable; all fields of characteristic $0$ and all finite fields
    are perfect.
\end{itemize}


%% ══════════════════════════════════════════════════════════
%% CHAPTER 10: ALGEBRAIC CLOSURE
%% ══════════════════════════════════════════════════════════
\chapter{Algebraic Closure}
\label{ch:algebraic-closure}
\index{algebraic closure}

\section{Algebraically closed fields}
\label{sec:alg-closed}

\begin{definition}[Algebraically closed field]
\label{def:alg-closed}
\index{algebraically closed field}
A field $K$ is \textbf{algebraically closed} if every nonconstant
polynomial in $K[X]$ has a root in $K$, or equivalently, every
nonconstant polynomial in $K[X]$ splits completely into linear factors
over $K$.
\end{definition}

\begin{proposition}[Equivalent characterisations]
\label{prop:alg-closed-equiv}
For a field $K$, the following are equivalent:
\begin{enumerate}[label=(\roman*)]
  \item $K$ is algebraically closed.
  \item Every nonconstant $f \in K[X]$ splits into linear factors.
  \item The only irreducible polynomials in $K[X]$ are the linear ones.
  \item $K$ has no proper algebraic extension.
\end{enumerate}
\end{proposition}

\begin{proof}
(i)$\Leftrightarrow$(ii): If every nonconstant polynomial has a root
$\alpha \in K$, we can factor out $(X-\alpha)$ and repeat by
induction on degree to split $f$ completely. The converse is
immediate.

(ii)$\Rightarrow$(iii): If $f \in K[X]$ is irreducible, then $f$
splits into linear factors, so $\deg f = 1$.

(iii)$\Rightarrow$(iv): Let $L/K$ be algebraic and $\alpha \in L$.
Then $\irr(\alpha, K)$ is irreducible, hence linear by (iii), so
$\alpha \in K$. Thus $L = K$.

(iv)$\Rightarrow$(i): Let $f \in K[X]$ be nonconstant with
irreducible factor $p$. Then $K[X]/(p)$ is an algebraic extension
of $K$, so $K[X]/(p) = K$ by (iv), giving $\deg p = 1$. Hence $f$
has a root in $K$.
\end{proof}

\begin{definition}[Algebraic closure]
\label{def:alg-closure}
\index{algebraic closure!definition}
An \textbf{algebraic closure} of a field $K$ is a field
$\overline{K} \supseteq K$ such that:
\begin{enumerate}[label=(\roman*)]
  \item $\overline{K}$ is algebraically closed.
  \item $\overline{K}/K$ is algebraic.
\end{enumerate}
\end{definition}

\begin{remark}[Minimality]
\label{rem:alg-closure-minimal}
Condition~(ii) means that $\overline{K}$ is the ``smallest''
algebraically closed field containing $K$. Any algebraically closed
field $\Omega \supseteq K$ contains (an isomorphic copy of)
$\overline{K}$ as a subfield.
\end{remark}

\section{Existence and uniqueness}
\label{sec:alg-closure-existence}

\begin{theorem}[Existence of algebraic closure]
\label{thm:alg-closure-existence}
\index{algebraic closure!existence}
Every field $K$ has an algebraic closure.
\end{theorem}

\begin{proof}[Proof sketch]
We outline the classical proof using Zorn's lemma\index{Zorn's lemma}.
Consider a ``sufficiently large'' set $\Omega \supseteq K$ and let
$\mathcal{F}$ be the collection of all algebraic extensions $L/K$
with $L \subseteq \Omega$, partially ordered by inclusion.

Every chain in $\mathcal{F}$ has an upper bound (the union of the
chain is again an algebraic extension of $K$). By Zorn's lemma,
$\mathcal{F}$ has a maximal element $\overline{K}$.

We claim $\overline{K}$ is algebraically closed. If not, there exists
a nonconstant irreducible $f \in \overline{K}[X]$ with $\deg f \ge 2$.
Then $\overline{K}[X]/(f)$ is a proper algebraic extension of
$\overline{K}$, and since $\overline{K}/K$ is algebraic, the
extension $\overline{K}[X]/(f)$ is also algebraic over $K$ (by
Theorem~\ref{thm:algebraic-tower}). One can embed
$\overline{K}[X]/(f)$ into $\Omega$ (with care about set-theoretic
issues), contradicting the maximality of $\overline{K}$.

The main technical subtlety is the choice of $\Omega$ to avoid
set-theoretic paradoxes. Artin's original proof handles this by
constructing a single polynomial ring in sufficiently many variables
and quotienting by an appropriate ideal. We refer the reader to
Lang's \textit{Algebra} for the complete argument.
\end{proof}

\begin{theorem}[Uniqueness of algebraic closure]
\label{thm:alg-closure-uniqueness}
\index{algebraic closure!uniqueness}
Any two algebraic closures of $K$ are $K$-isomorphic. The isomorphism
is not unique in general.
\end{theorem}

\begin{proof}[Proof sketch]
The proof uses the same extension-of-isomorphisms technique as
Theorem~\ref{thm:splitting-uniqueness}, together with Zorn's lemma.
One considers the set of all pairs $(L, \sigma)$ where $K \subseteq L
\subseteq \overline{K}$ and $\sigma \colon L \hookrightarrow
\overline{K}'$ is a $K$-embedding. Zorn's lemma gives a maximal such
pair, and maximality forces $L = \overline{K}$ and $\sigma$ to be
surjective.
\end{proof}

\section{The Fundamental Theorem of Algebra}
\label{sec:FTA}
\index{Fundamental Theorem of Algebra}

\begin{theorem}[Fundamental Theorem of Algebra]
\label{thm:FTA}
The field $\C$ of complex numbers is algebraically closed.
Equivalently, $\C$ is an algebraic closure of $\R$.
\end{theorem}

\begin{remark}[Historical note]
\label{rem:FTA-history}
The theorem was first stated by d'Alembert (1746)\index{d'Alembert}
and proved with varying degrees of rigour by Euler, Lagrange, Laplace,
and others. The first complete proof is generally attributed to
Gauss\index{Gauss} (1799), who gave four proofs during his lifetime.
Modern proofs use complex analysis (Liouville's theorem), topology
(winding numbers), or algebra combined with the intermediate value
theorem. We will return to give an algebraic proof using Galois theory
after developing the necessary machinery.
\end{remark}

\section{Finite fields}
\label{sec:finite-fields}

We now develop the complete theory of finite fields, one of the most
beautiful applications of the ideas developed so far.

\begin{theorem}[Structure of finite fields]
\label{thm:finite-field-structure}
\index{finite field!structure}
\begin{enumerate}[label=(\roman*)]
  \item Every finite field has order $p^n$ for some prime $p$ and
    integer $n \ge 1$.
  \item For each prime $p$ and integer $n \ge 1$, there exists a
    field of order $p^n$, unique up to isomorphism. We denote it
    $\F_{p^n}$.
  \item $\F_{p^n}$ is a splitting field of $X^{p^n} - X$ over $\F_p$.
\end{enumerate}
\end{theorem}

\begin{proof}
\textbf{(i)} Let $K$ be a finite field with $\chr(K) = p$. Then the
prime subfield of $K$ is $\F_p$, and $K$ is a finite-dimensional
$\F_p$-vector space. If $[K:\F_p] = n$, then $|K| = p^n$.

\textbf{(iii) and existence in (ii).} We show that $\F_{p^n}$ can be
realised as the splitting field of $f(X) = X^{p^n} - X$ over $\F_p$.

First, $f'(X) = p^n X^{p^n - 1} - 1 = -1$ (since $p^n \equiv 0$ in
$\F_p$), so $\gcd(f, f') = 1$, and $f$ has $p^n$ distinct roots in
any splitting field $L$ of $f$ over $\F_p$.

Let $S = \{ \alpha \in L : \alpha^{p^n} = \alpha \}$ be the set of
roots of $f$ in $L$. We claim $S$ is a subfield of $L$ with
$|S| = p^n$.

\textbf{$S$ is a subfield:}
\begin{itemize}
  \item $0^{p^n} = 0$ and $1^{p^n} = 1$, so $0, 1 \in S$.
  \item If $\alpha, \beta \in S$, then
    $(\alpha - \beta)^{p^n} = \alpha^{p^n} - \beta^{p^n}
    = \alpha - \beta$ (using the Frobenius), so
    $\alpha - \beta \in S$.
  \item If $\alpha, \beta \in S$ with $\beta \ne 0$, then
    $(\alpha\beta^{-1})^{p^n} = \alpha^{p^n}(\beta^{p^n})^{-1}
    = \alpha \beta^{-1}$, so $\alpha\beta^{-1} \in S$.
\end{itemize}
Since $f$ has exactly $p^n$ distinct roots, $|S| = p^n$.

Since $S$ is a subfield of $L$ containing all roots of $f$ and
$\F_p \subseteq S$, and $L$ is the splitting field of $f$ over
$\F_p$ (hence the smallest field containing $\F_p$ and all roots of
$f$), we conclude $L = S$. Thus $|L| = p^n$.

\textbf{Uniqueness in (ii).} Suppose $K$ is any field with
$|K| = p^n$. Then $K^{\times}$ is a group of order $p^n - 1$, so
$\alpha^{p^n - 1} = 1$ for all $\alpha \in K^{\times}$, hence
$\alpha^{p^n} = \alpha$ for all $\alpha \in K^{\times}$. This also
holds for $\alpha = 0$. Therefore every element of $K$ is a root of
$X^{p^n} - X$, and $K$ is a splitting field of $X^{p^n} - X$ over
$\F_p$. By the uniqueness of splitting fields
(Theorem~\ref{thm:splitting-uniqueness}), $K \cong \F_{p^n}$.
\end{proof}

\subsection{The subfield lattice}
\label{subsec:subfield-lattice}
\index{subfield lattice}

\begin{theorem}[Subfields of $\F_{p^n}$]
\label{thm:subfields-Fpn}
\index{finite field!subfields}
The subfields of $\F_{p^n}$ are precisely the fields $\F_{p^d}$ where
$d \mid n$. For each such divisor $d$ there is a unique subfield
isomorphic to $\F_{p^d}$.
\end{theorem}

\begin{proof}
If $\F_{p^d} \subseteq \F_{p^n}$, then $\F_{p^n}$ is a vector space
over $\F_{p^d}$, so $p^n = (p^d)^{[F_{p^n}:F_{p^d}]}$, giving
$n = d \cdot [\F_{p^n}:\F_{p^d}]$, hence $d \mid n$.

Conversely, suppose $d \mid n$. Then $p^d - 1 \mid p^n - 1$ (since
$X^d - 1 \mid X^n - 1$ in $\Z[X]$ when $d \mid n$), so
$X^{p^d} - X \mid X^{p^n} - X$ in $\F_p[X]$. (To verify this
divisibility: $\alpha^{p^d} = \alpha$ implies $\alpha^{p^n} = \alpha$,
which follows by applying the Frobenius $\varphi^d$ repeatedly:
$\alpha^{p^{2d}} = (\alpha^{p^d})^{p^d} = \alpha^{p^d} = \alpha$,
and by induction $\alpha^{p^{kd}} = \alpha$ for all $k$; taking
$k = n/d$ gives $\alpha^{p^n} = \alpha$.)

Hence every root of $X^{p^d} - X$ in $\F_{p^n}$ is also a root of
$X^{p^n} - X$. Since $X^{p^d} - X$ has exactly $p^d$ roots (it is
separable, as shown in the proof of
Theorem~\ref{thm:finite-field-structure}), the set
$S_d = \{\alpha \in \F_{p^n} : \alpha^{p^d} = \alpha\}$ has exactly
$p^d$ elements and is a subfield (by the same argument as before).
This $S_d$ is the unique subfield of $\F_{p^n}$ isomorphic to
$\F_{p^d}$.
\end{proof}

\begin{example}[Subfield lattice of $\F_{p^{12}}$]
\label{ex:subfield-lattice-p12}
\index{subfield lattice!of $\F_{p^{12}}$}
The divisors of $12$ are $1, 2, 3, 4, 6, 12$, and the divisibility
relations among them give the following lattice:
\end{example}

\begin{center}
\begin{tikzpicture}[
  node distance=1.8cm and 2.2cm,
  every node/.style={font=\small}
]
  \node (F1)  {$\F_p$};
  \node (F2)  [above left=of F1]  {$\F_{p^2}$};
  \node (F3)  [above right=of F1] {$\F_{p^3}$};
  \node (F4)  [above=of F2]       {$\F_{p^4}$};
  \node (F6)  [above=of F3]       {$\F_{p^6}$};
  \node (F12) [above right=of F4] {$\F_{p^{12}}$};

  \draw (F1)  -- (F2);
  \draw (F1)  -- (F3);
  \draw (F2)  -- (F4);
  \draw (F2)  -- (F6);
  \draw (F3)  -- (F6);
  \draw (F4)  -- (F12);
  \draw (F6)  -- (F12);
\end{tikzpicture}
\end{center}

Each line segment represents a field extension; the lower field is a
subfield of the upper field.

\subsection{The multiplicative group of a finite field}
\label{subsec:mult-group-cyclic}

\begin{theorem}[Multiplicative group is cyclic]
\label{thm:finite-field-mult-cyclic}
\index{finite field!multiplicative group}
\index{cyclic group!multiplicative group of a finite field}
For any finite field $\F_q$, the multiplicative group $\F_q^{\times}$
is cyclic of order $q - 1$.
\end{theorem}

\begin{proof}
Let $q - 1 = p_1^{a_1} p_2^{a_2} \cdots p_r^{a_r}$ be the prime
factorisation of $q - 1$. We will show that $\F_q^{\times}$ contains
an element of order $p_i^{a_i}$ for each $i$; their product then has
order $q - 1$, proving cyclicity.

Fix $i$ and write $q - 1 = p_i^{a_i} \cdot m$ where
$m = (q-1)/p_i^{a_i}$.

\textbf{Step 1.} The polynomial $X^m - 1 \in \F_q[X]$ has degree $m$
and therefore at most $m$ roots in $\F_q$. Since
$|\F_q^{\times}| = q - 1 > m$, there exists an element
$b \in \F_q^{\times}$ with $b^m \ne 1$.

\textbf{Step 2.} Set $c = b^{(q-1)/p_i^{a_i}} = b^m$. Wait ---
let us be more careful. We want an element of order exactly
$p_i^{a_i}$. Consider $\alpha = b^{m/1} = b^{(q-1)/p_i^{a_i}}$.
No, let us use a cleaner argument.

Set $\beta = b^{(q-1)/p_i^{a_i}}$. Then
$\beta^{p_i^{a_i}} = b^{q-1} = 1$, so $\mathrm{ord}(\beta) \mid
p_i^{a_i}$, hence $\mathrm{ord}(\beta) = p_i^{e_i}$ for some
$0 \le e_i \le a_i$.

\textbf{Step 3.} We claim $e_i \ge 1$. If $e_i = 0$, then
$\beta = 1$, so $b^{(q-1)/p_i^{a_i}} = 1$, i.e.\ $b^m = 1$,
contradicting the choice of $b$. So $e_i \ge 1$.

\textbf{Step 4.} Now set $\gamma_i = \beta^{p_i^{e_i - a_i}}$. Wait
--- since $e_i \le a_i$, set
$\gamma_i = \beta^{p_i^{e_i}/p_i^{a_i}}$, but this need not be an
integer. Let us instead use the following cleaner approach.

Since $\mathrm{ord}(\beta) = p_i^{e_i}$ with $e_i \ge 1$, the
element $\gamma_i = \beta^{p_i^{e_i - a_i}}$ if $e_i = a_i$ gives
$\gamma_i = \beta$ with order $p_i^{a_i}$, and we are done.

If $e_i < a_i$, we need to modify our choice. Consider the set
$S = \{ x^m : x \in \F_q^{\times} \}$. This is a subgroup of
$\F_q^{\times}$. An element $x \in \F_q^{\times}$ satisfies
$x^m = 1$ if and only if $x$ is a root of $X^m - 1$, and there are
at most $m$ such elements. By the homomorphism
$\F_q^{\times} \to \F_q^{\times}$, $x \mapsto x^m$, we have
$|S| = |\F_q^{\times}| / |\Ker| = (q-1)/|\{x : x^m = 1\}|$.

Since $X^m - 1$ divides $X^{q-1} - 1$ in $\F_q[X]$ (because
$m \mid q-1$), and $X^{q-1} - 1$ has exactly $q-1$ roots in
$\F_q$, the polynomial $X^m - 1$ has exactly $m$ roots in $\F_q$.
(In a finite field, $X^d - 1$ has exactly $d$ roots whenever
$d \mid q-1$, which follows because $\F_q^{\times}$ is a group of
order $q-1$ and $X^{q-1} - 1 = \prod_{\alpha \in \F_q^{\times}}
(X - \alpha)$, so $X^d - 1 = \gcd(X^d - 1, X^{q-1}-1)$ has exactly
$d$ roots.)

Therefore $|S| = (q-1)/m = p_i^{a_i}$. So $S$ is a subgroup of
$\F_q^{\times}$ of order $p_i^{a_i}$.

\textbf{Step 5.} By the classification of finite abelian groups,
every subgroup of $\F_q^{\times}$ of order $p_i^{a_i}$ is a
$p_i$-group. We claim it is cyclic. Indeed, if it were not cyclic,
then every element $s \in S$ would satisfy
$s^{p_i^{a_i - 1}} = 1$, giving $|S|$ roots of the polynomial
$X^{p_i^{a_i -1}} - 1$, namely $p_i^{a_i}$ roots. But this
polynomial has degree $p_i^{a_i - 1} < p_i^{a_i}$, so it can have
at most $p_i^{a_i-1}$ roots in the field $\F_q$ --- a
contradiction.

Therefore $S$ contains an element $\gamma_i$ of order $p_i^{a_i}$.

\textbf{Step 6.} Having found $\gamma_i \in \F_q^{\times}$ of order
$p_i^{a_i}$ for each $1 \le i \le r$, set
$\gamma = \gamma_1 \gamma_2 \cdots \gamma_r$.

Since $\gcd(p_i^{a_i}, p_j^{a_j}) = 1$ for $i \ne j$, and the
$\gamma_i$ are elements of the abelian group $\F_q^{\times}$, we have
\[
  \mathrm{ord}(\gamma) =
  \mathrm{lcm}(\mathrm{ord}(\gamma_1), \dots,
  \mathrm{ord}(\gamma_r)).
\]
Actually, we need to be slightly more careful. Since each $\gamma_i$
has order $p_i^{a_i}$ and $\gamma_i \in S_i = \{x^{m_i} : x \in
\F_q^{\times}\}$ where $m_i = (q-1)/p_i^{a_i}$, the subgroups
$S_i$ and $S_j$ (for $i \ne j$) satisfy $S_i \cap S_j = \{1\}$
(since $|S_i \cap S_j|$ divides $\gcd(|S_i|, |S_j|) = 1$). It
follows that the product $\gamma = \gamma_1 \cdots \gamma_r$ has
order
\[
  \mathrm{ord}(\gamma) = \prod_{i=1}^r p_i^{a_i} = q - 1.
\]
Therefore $\F_q^{\times} = \langle \gamma \rangle$ is cyclic.
\end{proof}

\begin{remark}[Generators]
\label{rem:primitive-element-finite}
\index{primitive element!of a finite field}
A generator of $\F_q^{\times}$ is called a \textbf{primitive element}
of $\F_q$. By the theory of cyclic groups, there are exactly
$\varphi(q-1)$ primitive elements.
\end{remark}

\begin{corollary}[Finite subgroups of multiplicative groups]
\label{cor:finite-subgroup-mult}
Let $K$ be any field and let $G$ be a finite subgroup of $K^{\times}$.
Then $G$ is cyclic.
\end{corollary}

\begin{proof}
The proof of Step~5 above applies verbatim: if $|G| = n$ and $G$ were
not cyclic, then every element of $G$ would satisfy $x^{n/p} = 1$ for
some prime $p \mid n$, but $X^{n/p} - 1$ has at most $n/p < n$ roots
in the field $K$.

More precisely, let $|G| = n$ and suppose $G$ is not cyclic. Then the
exponent of $G$ (the least common multiple of the orders of all
elements) is $e < n$. Every $g \in G$ satisfies $g^e = 1$, so $G$
consists of roots of $X^e - 1$. But this polynomial has at most $e$
roots in the field $K$, giving $n = |G| \le e < n$, a contradiction.
\end{proof}

\section{Exercises}
\label{sec:ch10-exercises}

\begin{exercise}[Algebraically closed $\Rightarrow$ infinite]
\label{exer:alg-closed-infinite}
Prove that every algebraically closed field is infinite.
\textit{Hint:} Adapt Euclid's proof that there are infinitely many
primes.
\end{exercise}

\begin{exercise}[Algebraic closure of $\F_p$]
\label{exer:alg-closure-Fp}
\index{algebraic closure!of $\F_p$}
Show that $\overline{\F_p} = \bigcup_{n \ge 1} \F_{p^n}$ (where the
union is taken inside a fixed algebraic closure). Verify that this
union is indeed a field and that it is algebraically closed.
\end{exercise}

\begin{exercise}[Subfields of $\F_{p^{30}}$]
\label{exer:subfields-Fp30}
List all subfields of $\F_{p^{30}}$ and draw the inclusion lattice.
\end{exercise}

\begin{exercise}[Number of irreducible polynomials]
\label{exer:number-irred-polys}
\index{irreducible polynomial!counting over finite fields}
Let $N_q(n)$ denote the number of monic irreducible polynomials of
degree $n$ over $\F_q$. Show that
\[
  q^n = \sum_{d \mid n} d\, N_q(d),
\]
and deduce by M\"obius inversion that
\[
  N_q(n) = \frac{1}{n} \sum_{d \mid n} \mu(n/d)\, q^d,
\]
where $\mu$ is the M\"obius function.
\end{exercise}

\begin{exercise}[Explicit irreducible over $\F_2$]
\label{exer:explicit-irred-F2}
Find all irreducible polynomials of degree $4$ over $\F_2$. How many
are there? Verify your answer using the formula from
Exercise~\ref{exer:number-irred-polys}.
\end{exercise}

\begin{exercise}[Frobenius generates the Galois group]
\label{exer:frobenius-galois}
Show that $\Gal(\F_{p^n}/\F_p)$ is cyclic of order $n$, generated
by the Frobenius automorphism $\varphi \colon x \mapsto x^p$.
\end{exercise}

\begin{exercise}[Finite field extensions]
\label{exer:ff-extensions}
Show that $\F_{p^m} \subseteq \F_{p^n}$ if and only if $m \mid n$.
When $m \mid n$, show $[\F_{p^n}:\F_{p^m}] = n/m$.
\end{exercise}

\begin{exercise}[Primitive elements]
\label{exer:primitive-elements}
\begin{enumerate}[label=(\alph*)]
  \item Find a primitive element (generator of $\F_q^{\times}$) for
    $\F_7$, $\F_{11}$, and $\F_{13}$.
  \item Find a primitive element for $\F_4$ and $\F_8$.
  \item Show that $2$ is a primitive element of $\F_{11}$ but not of
    $\F_{13}$.
\end{enumerate}
\end{exercise}

\begin{exercise}[Wedderburn's little theorem (guided)]
\label{exer:wedderburn}
\index{Wedderburn's little theorem}
Let $D$ be a finite division ring (not assumed commutative).
\begin{enumerate}[label=(\alph*)]
  \item Show that the centre $Z(D)$ is a field, say $Z(D) \cong
    \F_q$ with $q = p^m$.
  \item Show $|D| = q^n$ for some $n \ge 1$.
  \item For each $x \in D$, the centraliser $C_D(x)$ is a division
    subring containing $Z(D)$, so $|C_D(x)| = q^{d(x)}$ with
    $d(x) \mid n$.
  \item Write the class equation for $D^{\times}$ and use cyclotomic
    polynomial arguments to derive a contradiction if $n > 1$.
\end{enumerate}
This proves Wedderburn's theorem: every finite division ring is a
field.
\end{exercise}

\begin{exercise}[Artin--Schreier polynomials]
\label{exer:artin-schreier}
\index{Artin--Schreier polynomial}
Let $K$ be a field of characteristic $p > 0$ and $a \in K$ such that
$X^p - X - a$ has no root in $K$. Show that $X^p - X - a$ is
irreducible over $K$.
\textit{Hint:} If $\alpha$ is a root, show that all roots are
$\alpha, \alpha+1, \dots, \alpha + (p-1)$.
\end{exercise}

\subsection*{Chapter summary}
\addcontentsline{toc}{section}{Chapter summary}

\begin{itemize}[leftmargin=2em]
  \item A field is \textbf{algebraically closed} if every nonconstant
    polynomial has a root; equivalently, the only irreducible
    polynomials are linear.
  \item Every field $K$ has an \textbf{algebraic closure}
    $\overline{K}$, unique up to $K$-isomorphism (existence uses
    Zorn's lemma).
  \item The \textbf{Fundamental Theorem of Algebra} asserts that $\C$
    is algebraically closed.
  \item For each prime power $q = p^n$, there is a unique (up to
    isomorphism) finite field $\F_q$, which is the splitting field of
    $X^q - X$ over $\F_p$.
  \item The subfields of $\F_{p^n}$ are exactly $\F_{p^d}$ for
    $d \mid n$.
  \item The multiplicative group $\F_q^{\times}$ is \textbf{cyclic} of
    order $q - 1$.
\end{itemize}
%%%%%%%%%%%%%%%%%%%%%%%%%%%%%%%%%%%%%%%%%%%%%%%%%%%%%%%%%%%%
%% ABSTRACT ALGEBRA II — PART 2 (Chapters 11–13)
%% English Edition
%%%%%%%%%%%%%%%%%%%%%%%%%%%%%%%%%%%%%%%%%%%%%%%%%%%%%%%%%%%%

%% ============================================================
%%  CHAPTER 11 — GALOIS EXTENSIONS
%% ============================================================
\chapter{Galois Extensions}\label{ch:galois-extensions}
\index{Galois extension}

In the preceding chapters we studied the notions of normality and
separability for algebraic field extensions.  We now bring them together
in the central concept of \emph{Galois extension} and begin the study of
the group of automorphisms that will lead, in Chapter~\ref{ch:fundamental-theorem},
to the Fundamental Theorem of Galois Theory.

%% ------------------------------------------------------------
\section{The Automorphism Group and Fixed Fields}
\label{sec:aut-fixed}
\index{automorphism group}\index{fixed field}

\begin{definition}[Automorphism group of an extension]\label{def:aut-extension}
Let $L/K$ be a field extension.  The \emph{automorphism group of $L/K$} is
\[
  \Aut(L/K) \;=\; \bigl\{\,\sigma\in\Aut(L) \;\bigm|\;
        \sigma(a)=a \text{ for all } a\in K\bigr\}.
\]
When $L/K$ is a Galois extension (Definition~\ref{def:galois-ext}), we write
$\Gal(L/K)$\index{Galois group} instead of $\Aut(L/K)$ and call it the
\emph{Galois group}\index{Galois group} of the extension.
\end{definition}

\begin{definition}[Fixed field]\label{def:fixed-field}
Let $L$ be a field and $H\leq\Aut(L)$ a subgroup.  The \emph{fixed field}
of $H$ is
\[
  \Fix(H) \;=\; L^{H} \;=\;
  \bigl\{\,a\in L \;\bigm|\; \sigma(a)=a \text{ for all } \sigma\in H\bigr\}.
\]
\end{definition}

\begin{remark}[Elementary properties]\label{rem:fixed-elementary}
\begin{enumerate}
\item $L^{H}$ is indeed a subfield of $L$.
\item If $H_{1}\leq H_{2}\leq\Aut(L)$, then $L^{H_{2}}\subseteq L^{H_{1}}$.
\item For any subgroup $H\leq\Aut(L)$, one has
      $H\leq\Aut\!\bigl(L/L^{H}\bigr)$.
\end{enumerate}
\end{remark}

%% ------------------------------------------------------------
\section{Galois Extensions: Definition and Characterisations}
\label{sec:galois-def}

\begin{definition}[Galois extension]\label{def:galois-ext}
A finite extension $L/K$ is called a \emph{Galois extension}\index{Galois extension}
if it is both \emph{normal} and \emph{separable}.
\end{definition}

\begin{theorem}[Equivalent characterisations of Galois extensions]
\label{thm:galois-equiv}
Let $L/K$ be a finite extension of degree $n=[L:K]$.  The following are
equivalent:
\begin{enumerate}
\item\label{it:gal-ns} $L/K$ is normal and separable.
\item\label{it:gal-split} $L$ is the splitting field of some separable
      polynomial $f\in K[X]$.
\item\label{it:gal-aut} $|\!\Aut(L/K)|=n$.
\item\label{it:gal-fix} $L^{\Aut(L/K)}=K$ (the fixed field of $\Aut(L/K)$ is $K$).
\end{enumerate}
\end{theorem}

We prove the equivalences in several steps across this section.

%% --- NORMAL ⟺ SPLITTING FIELD ---
\subsection{Normal Extensions and Splitting Fields}
\label{subsec:normal-splitting}
\index{normal extension}\index{splitting field}

\begin{theorem}[Normal $\Longleftrightarrow$ splitting field]\label{thm:normal-splitting}
Let $L/K$ be a finite extension.  Then $L/K$ is normal if and only if
$L$ is the splitting field over $K$ of some polynomial $f\in K[X]$.
\end{theorem}

\begin{proof}
$(\Longrightarrow)$\;
Suppose $L/K$ is normal and write $L=K(\alpha_{1},\dots,\alpha_{r})$
for finitely many algebraic elements $\alpha_{i}$.  Let
$f_{i}=\irr(\alpha_{i},K)\in K[X]$ be the minimal polynomial of $\alpha_{i}$
over $K$.  By normality, every irreducible polynomial in $K[X]$ having a root
in $L$ splits completely in $L$; in particular each $f_{i}$ splits in $L[X]$.
Set $f=f_{1}f_{2}\cdots f_{r}$.  Then $f\in K[X]$, all roots of $f$ lie
in $L$, and $L$ is generated over $K$ by these roots.  Hence $L$ is a
splitting field of $f$ over $K$.

$(\Longleftarrow)$\;
Suppose $L$ is the splitting field of $f\in K[X]$.  We must show that
every irreducible $g\in K[X]$ having a root $\beta\in L$ splits
completely in $L$.  Let $\beta'\in\overline{K}$ be any root of $g$.
Since $g$ is irreducible over $K$, the map $\beta\mapsto\beta'$ extends
to a $K$-isomorphism $\varphi\colon K(\beta)\xrightarrow{\;\sim\;}K(\beta')$.
Because $L$ is a splitting field of $f$ over $K$, and $K(\beta)\subseteq L$,
we may extend $\varphi$ to an embedding $\tilde{\varphi}\colon L\hookrightarrow\overline{K}$.
But $L$, being a splitting field, is mapped to itself by any such embedding
(the roots of $f$ are permuted), so $\tilde{\varphi}(L)=L$.  In particular
$\beta'=\tilde{\varphi}(\beta)\in L$.  Since $\beta'$ was an arbitrary root
of $g$, the polynomial $g$ splits in $L$.
\end{proof}

%% --- |Gal(L/K)| = [L:K] ---
\subsection{Order of the Galois Group}
\label{subsec:order-galois}

\begin{theorem}[{$|\Gal(L/K)|=[L:K]$ for Galois extensions}]
\label{thm:gal-order}
\index{Galois group!order}
If $L/K$ is a finite Galois extension, then
$|\Gal(L/K)|=[L:K]$.
\end{theorem}

\begin{proof}
Set $n=[L:K]$.  We prove both inequalities.

\medskip\noindent
\textbf{Step 1:} $|\!\Aut(L/K)|\leq n$.

Write $L=K(\alpha_{1},\dots,\alpha_{r})$.  Every $\sigma\in\Aut(L/K)$ is
determined by the images $\sigma(\alpha_{1}),\dots,\sigma(\alpha_{r})$.
For each $i$, the element $\sigma(\alpha_{i})$ must be a root of
$\irr(\alpha_{i},K)$, so there are at most $\deg\irr(\alpha_{i},K)$
choices.  A standard induction on $r$ using the tower
\[
  K\;\subseteq\; K(\alpha_{1})\;\subseteq\;
  K(\alpha_{1},\alpha_{2})\;\subseteq\;\cdots\;\subseteq\; L
\]
shows that the number of $K$-embeddings of $L$ into $\overline{K}$ is at most
\[
  [K(\alpha_{1}):K]\cdot[K(\alpha_{1},\alpha_{2}):K(\alpha_{1})]
  \cdots = [L:K] = n.
\]
Since every element of $\Aut(L/K)$ restricts to such a $K$-embedding, we get
$|\!\Aut(L/K)|\leq n$.

\medskip\noindent
\textbf{Step 2:} $|\!\Aut(L/K)|\geq n$.

Since $L/K$ is separable, we may write $L=K(\alpha)$ by the Primitive
Element Theorem (Theorem~\ref{thm:primitive-element}; we prove it
independently in Chapter~\ref{ch:fundamental-theorem}).  Let
$f=\irr(\alpha,K)$.  Then $\deg f=[K(\alpha):K]=n$.  Because $L/K$ is
normal, $f$ splits in $L$; because $L/K$ is separable, $f$ has $n$ distinct
roots $\alpha_{1},\dots,\alpha_{n}\in L$.  For each $i$ the assignment
$\alpha\mapsto\alpha_{i}$ extends to a unique $K$-automorphism
$\sigma_{i}\in\Aut(L/K)$.  These $\sigma_{i}$ are pairwise distinct,
so $|\!\Aut(L/K)|\geq n$.

\medskip
Combining Steps 1 and 2 yields $|\!\Aut(L/K)|=n=[L:K]$.
\end{proof}

%% --- ARTIN'S LEMMA ---
\subsection{Artin's Lemma}
\label{subsec:artin-lemma}

\begin{lemma}[Artin's lemma]\label{lem:artin}
\index{Artin's lemma}
Let $L$ be a field and $G$ a finite subgroup of $\Aut(L)$ with $|G|=n$.
Then $[L:L^{G}]=n$.  In particular $L/L^{G}$ is a Galois extension with
$\Gal(L/L^{G})=G$.
\end{lemma}

\begin{proof}
Set $K=L^{G}$.  We prove $[L:K]=n$.

\medskip\noindent
\textbf{Step 1:} $[L:K]\leq n$.

Suppose for contradiction that $[L:K]>n$.  Choose $n+1$ elements
$\alpha_{1},\dots,\alpha_{n+1}\in L$ that are $K$-linearly independent.
Write $G=\{\sigma_{1},\dots,\sigma_{n}\}$.  Consider the homogeneous
system of $n$ linear equations in $n+1$ unknowns $x_{1},\dots,x_{n+1}$:
\[
  \sum_{j=1}^{n+1}\sigma_{i}(\alpha_{j})\,x_{j}=0,
  \qquad i=1,\dots,n.
\]
Since there are more unknowns than equations, there exists a non-trivial
solution $(c_{1},\dots,c_{n+1})\in L^{n+1}\setminus\{0\}$.  Choose such a
solution with the fewest non-zero entries.  After reordering, assume
$c_{1},\dots,c_{m}\neq 0$ and $c_{m+1}=\cdots=c_{n+1}=0$, with $m$
minimal.  Dividing by $c_{m}$, we may assume $c_{m}=1$.

If all $c_{j}\in K$, then taking $\sigma_{i}=\mathrm{id}$ gives
$\sum_{j}c_{j}\alpha_{j}=0$, contradicting the $K$-linear independence of
the $\alpha_{j}$.  So some $c_{j}\notin K$; say $c_{1}\notin K$.  Then
there exists $\tau\in G$ with $\tau(c_{1})\neq c_{1}$.  Applying $\tau$ to
the system:
\[
  \sum_{j=1}^{m}\tau(\sigma_{i}(\alpha_{j}))\,\tau(c_{j})=0
  \quad\text{for all }i.
\]
Since $\tau$ permutes $G$, $\{\tau\sigma_{1},\dots,\tau\sigma_{n}\}=G$, so
this is the same system with $c_{j}$ replaced by $\tau(c_{j})$.
Subtracting from the original:
\[
  \sum_{j=1}^{m}\sigma_{i}(\alpha_{j})\bigl(c_{j}-\tau(c_{j})\bigr)=0
  \quad\text{for all }i.
\]
Since $c_{m}=\tau(c_{m})=1$, the $m$-th coefficient vanishes, and
$c_{1}-\tau(c_{1})\neq 0$.  This gives a non-trivial solution with fewer
than $m$ non-zero entries, contradicting the minimality of $m$.

\medskip\noindent
\textbf{Step 2:} $[L:K]\geq n$.

We always have $G\leq\Aut(L/K)$.  By Step~1, $[L:K]\leq n$, and
Theorem~\ref{thm:gal-order} (Step 1 of its proof, which applies to any
finite extension) gives $|\!\Aut(L/K)|\leq[L:K]$.  Hence
$n=|G|\leq|\!\Aut(L/K)|\leq[L:K]\leq n$, so all inequalities are
equalities.

\medskip
Therefore $[L:K]=n$, $\Aut(L/K)=G$, and $L/K$ is Galois.
\end{proof}

%% --- Proof of the remaining equivalences ---
\begin{proof}[Proof of Theorem~\ref{thm:galois-equiv}]
$\ref{it:gal-ns}\Rightarrow\ref{it:gal-split}$:
By Theorem~\ref{thm:normal-splitting}, normality implies $L$ is a splitting
field of some $f\in K[X]$.  Since $L/K$ is separable, the roots of the
minimal polynomials generating $L$ are separable, and we can choose $f$
to be separable (take $f$ to be the product of the distinct irreducible
factors of the minimal polynomials).

$\ref{it:gal-split}\Rightarrow\ref{it:gal-ns}$:
A splitting field of a separable polynomial is normal
(Theorem~\ref{thm:normal-splitting}).  Since the roots of $f$ are separable
over $K$, every element of $L=K(\alpha_{1},\dots,\alpha_{r})$ is separable
over $K$.

$\ref{it:gal-ns}\Rightarrow\ref{it:gal-aut}$:
This is Theorem~\ref{thm:gal-order}.

$\ref{it:gal-aut}\Rightarrow\ref{it:gal-fix}$:
Set $G=\Aut(L/K)$ and $E=L^{G}$.  Then $K\subseteq E$ and
$G\leq\Aut(L/E)$.  Artin's lemma gives $[L:E]=|G|=n=[L:K]$, so $E=K$.

$\ref{it:gal-fix}\Rightarrow\ref{it:gal-ns}$:
Set $G=\Aut(L/K)$ and suppose $L^{G}=K$.  By Artin's lemma, $[L:K]=|G|$
and $L/K$ is Galois, hence normal and separable.
\end{proof}

%% ------------------------------------------------------------
\section{Detailed Examples of Galois Groups}
\label{sec:galois-examples}

\subsection{$\Gal(\Q(\sqrt{2}\,)/\Q)\cong\Z/2\Z$}
\label{subsec:ex-sqrt2}

\begin{example}[$\Gal(\Q(\sqrt{2}\,)/\Q)$]\label{ex:gal-sqrt2}
\index{Galois group!of $\Q(\sqrt{2}\,)/\Q$}
The polynomial $f=X^{2}-2$ is irreducible over $\Q$ (Eisenstein at $p=2$)
and separable.  Its splitting field is $\Q(\sqrt{2}\,)$, which is therefore
Galois over $\Q$ of degree $2$.

Any $\sigma\in\Gal(\Q(\sqrt{2}\,)/\Q)$ satisfies $\sigma(\sqrt{2}\,)^{2}=
\sigma(2)=2$, so $\sigma(\sqrt{2}\,)=\pm\sqrt{2}$.  This gives exactly
two automorphisms: the identity and $\sigma\colon\sqrt{2}\mapsto -\sqrt{2}$.
Hence
\[
  \Gal\!\bigl(\Q(\sqrt{2}\,)/\Q\bigr)
  \;=\;\{\mathrm{id},\,\sigma\}
  \;\cong\;\Z/2\Z.
\]
\end{example}

\subsection{$\Gal(\Q(\sqrt[3]{2},\omega)/\Q)\cong S_{3}$}
\label{subsec:ex-s3}

\begin{example}[{Full computation of $\Gal(\Q(\sqrt[3]{2},\omega)/\Q)\cong S_{3}$}]
\label{ex:gal-s3}
\index{Galois group!of $\Q(\sqrt[3]{2},\omega)/\Q$}
Let $\omega=e^{2\pi i/3}=(-1+\sqrt{-3}\,)/2$ be a primitive cube root of
unity and consider $f=X^{3}-2$, which is irreducible over $\Q$ by
Eisenstein at $p=2$.  Its three roots are
\[
  \alpha_{1}=\sqrt[3]{2},\qquad
  \alpha_{2}=\omega\sqrt[3]{2},\qquad
  \alpha_{3}=\omega^{2}\sqrt[3]{2}.
\]
The splitting field is $L=\Q(\sqrt[3]{2},\omega)$.

\medskip\noindent
\textbf{Degree computation.}\;
$[\Q(\sqrt[3]{2}):\Q]=3$ and $\omega\notin\Q(\sqrt[3]{2})\subseteq\R$, so
$[\Q(\sqrt[3]{2},\omega):\Q(\sqrt[3]{2})]=2$ (the minimal polynomial of
$\omega$ over $\Q(\sqrt[3]{2})$ is $X^{2}+X+1$).  Thus $[L:\Q]=6$.

\medskip\noindent
\textbf{The six automorphisms.}\;
By Theorem~\ref{thm:gal-order}, $|\Gal(L/\Q)|=6$.  Each $\sigma\in
\Gal(L/\Q)$ is determined by
\[
  \sigma(\sqrt[3]{2})\in\{\alpha_{1},\alpha_{2},\alpha_{3}\}
  \quad\text{and}\quad
  \sigma(\omega)\in\{\omega,\omega^{2}\}.
\]
Label the automorphisms by their action on the roots:
\[
\renewcommand{\arraystretch}{1.2}
\begin{array}{c|cc}
  & \sigma(\sqrt[3]{2}) & \sigma(\omega) \\\hline
  \mathrm{id} & \alpha_{1} & \omega \\
  \tau & \alpha_{1} & \omega^{2} \\
  \rho & \alpha_{2} & \omega \\
  \rho^{2} & \alpha_{3} & \omega \\
  \rho\tau & \alpha_{2} & \omega^{2} \\
  \rho^{2}\tau & \alpha_{3} & \omega^{2}
\end{array}
\]
Here $\rho$ acts on $\{\alpha_{1},\alpha_{2},\alpha_{3}\}$ as the $3$-cycle
$(1\;2\;3)$ and $\tau$ as the transposition induced by complex conjugation
(swapping $\omega\leftrightarrow\omega^{2}$, which acts as $(2\;3)$ on the
roots).  These generate $S_{3}$, confirming
\[
  \Gal\!\bigl(\Q(\sqrt[3]{2},\omega)/\Q\bigr)\;\cong\; S_{3}.
\]
\end{example}

\subsection{$\Gal(\F_{p^{n}}/\F_{p})\cong\Z/n\Z$ — The Frobenius Automorphism}
\label{subsec:ex-frobenius}

\begin{theorem}[Galois group of a finite field extension]\label{thm:gal-finite-field}
\index{Galois group!of finite fields}\index{Frobenius automorphism}
For every prime $p$ and integer $n\geq 1$, the extension $\F_{p^{n}}/\F_{p}$
is Galois with
\[
  \Gal(\F_{p^{n}}/\F_{p})
  \;=\;\langle\varphi\rangle
  \;\cong\;\Z/n\Z,
\]
where $\varphi\colon x\mapsto x^{p}$ is the \emph{Frobenius automorphism}.
\end{theorem}

\begin{proof}
\textbf{Step 1: $\varphi$ is a well-defined automorphism.}\;
For $x,y\in\F_{p^{n}}$ we have
\[
  \varphi(x+y)=(x+y)^{p}=x^{p}+y^{p}=\varphi(x)+\varphi(y)
\]
(the binomial coefficients $\binom{p}{k}$ with $0<k<p$ are divisible by $p$,
hence vanish in characteristic $p$), and
$\varphi(xy)=(xy)^{p}=x^{p}y^{p}=\varphi(x)\varphi(y)$.  Since $\varphi$
is a non-zero ring homomorphism from a field to itself, it is injective.
As $\F_{p^{n}}$ is finite, $\varphi$ is also surjective, hence an
automorphism.  For $a\in\F_{p}$, Fermat's little theorem gives
$\varphi(a)=a^{p}=a$, so $\varphi\in\Aut(\F_{p^{n}}/\F_{p})$.

\medskip\noindent
\textbf{Step 2: $\varphi$ has order exactly $n$.}\;
We have $\varphi^{k}(x)=x^{p^{k}}$.  Thus $\varphi^{k}=\mathrm{id}$ if
and only if $x^{p^{k}}=x$ for all $x\in\F_{p^{n}}$.  The elements of
$\F_{p^{n}}$ are precisely the roots of $X^{p^{n}}-X$, so $\varphi^{n}=
\mathrm{id}$.  If $\varphi^{k}=\mathrm{id}$ for some $k<n$, then every
element of $\F_{p^{n}}$ would be a root of $X^{p^{k}}-X$, which has at most
$p^{k}<p^{n}$ roots — a contradiction.

\medskip\noindent
\textbf{Step 3: Conclusion.}\;
We have $\langle\varphi\rangle\leq\Gal(\F_{p^{n}}/\F_{p})$ with
$|\langle\varphi\rangle|=n$.  Since $[\F_{p^{n}}:\F_{p}]=n$ and
$|\Gal(\F_{p^{n}}/\F_{p})|\leq n$, equality holds:
$\Gal(\F_{p^{n}}/\F_{p})=\langle\varphi\rangle\cong\Z/n\Z$.
\end{proof}

\subsection{Cyclotomic Extensions}
\label{subsec:cyclotomic}
\index{cyclotomic extension}\index{cyclotomic polynomial}

\begin{theorem}[Galois group of cyclotomic extensions]\label{thm:gal-cyclotomic}
Let $n\geq 1$ and $\zeta_{n}=e^{2\pi i/n}$.  Then $\Q(\zeta_{n})/\Q$ is
Galois with
\[
  \Gal\!\bigl(\Q(\zeta_{n})/\Q\bigr)
  \;\cong\; (\Z/n\Z)^{\times}.
\]
In particular $[\Q(\zeta_{n}):\Q]=\varphi(n)$, where $\varphi$ is Euler's
totient function.
\end{theorem}

\begin{proof}[Proof sketch]
The minimal polynomial of $\zeta_{n}$ over $\Q$ is the $n$-th cyclotomic
polynomial $\Phi_{n}\in\Z[X]$, which is irreducible over $\Q$ (a classical
result proved, e.g., by reduction modulo~$p$).  Hence
$[\Q(\zeta_{n}):\Q]=\deg\Phi_{n}=\varphi(n)$.  Since $\Q(\zeta_{n})$ is
the splitting field of $X^{n}-1$ (a separable polynomial), the extension is
Galois.  Every $\sigma\in\Gal(\Q(\zeta_{n})/\Q)$ satisfies
$\sigma(\zeta_{n})^{n}=1$, so $\sigma(\zeta_{n})=\zeta_{n}^{a}$ for a
unique $a\in(\Z/n\Z)^{\times}$.  The map $\sigma\mapsto a$ is readily
checked to be a group isomorphism.
\end{proof}

%% --- DIAGRAM ---
\begin{center}
\begin{tikzcd}[row sep=1.5em, column sep=1.5em]
  & \Q(\zeta_{n}) \arrow[dl,dash,"{\varphi(n)/\varphi(n/p)}"']
    \arrow[dr,dash,"{\varphi(n)/\varphi(n/q)}"] & \\
  \Q(\zeta_{n/p}) \arrow[dr,dash] & & \Q(\zeta_{n/q}) \arrow[dl,dash] \\
  & \Q &
\end{tikzcd}
\captionof{figure}{Subfield lattice of $\Q(\zeta_{n})$ for $n=pq$
(distinct primes).}\label{fig:cyclotomic-lattice}
\end{center}

%% ------------------------------------------------------------
\section{Exercises}
\label{sec:ch11-exercises}

\begin{exercise}[Automorphisms of $\Q(\sqrt{5}\,)$]\label{exo:aut-sqrt5}
Determine $\Gal(\Q(\sqrt{5}\,)/\Q)$ and all intermediate fields.
\end{exercise}

\begin{exercise}[Fourth roots of unity]\label{exo:fourth-roots}
Show that $\Q(i)/\Q$ is Galois, determine $\Gal(\Q(i)/\Q)$, and verify the
degree--order equality.
\end{exercise}

\begin{exercise}[Splitting field of $X^{4}-2$]\label{exo:x4-2}
Let $L$ be the splitting field of $X^{4}-2$ over $\Q$.
\begin{enumerate}
\item Show $[L:\Q]=8$.
\item Prove that $\Gal(L/\Q)\cong D_{4}$ (the dihedral group of order $8$).
\end{enumerate}
\end{exercise}

\begin{exercise}[Fixed field computation]\label{exo:fixed-field}
Let $L=\Q(\sqrt{2},\sqrt{3}\,)$ and $\sigma\colon\sqrt{2}\mapsto -\sqrt{2}$,
$\sqrt{3}\mapsto\sqrt{3}$.  Compute $L^{\langle\sigma\rangle}$.
\end{exercise}

\begin{exercise}[Artin's lemma in practice]\label{exo:artin-practice}
Let $\sigma\colon\C\to\C$ be complex conjugation and $G=\{1,\sigma\}$.
Use Artin's lemma to prove that $[\C:\R]=2$ without using any prior
knowledge of $\C/\R$.
\end{exercise}

\begin{exercise}[Frobenius powers]\label{exo:frobenius-powers}
Let $\varphi$ denote the Frobenius automorphism of $\F_{p^{6}}/\F_{p}$.
List all subgroups of $\langle\varphi\rangle$ and the corresponding fixed
fields.
\end{exercise}

\begin{exercise}[Cyclotomic polynomials]\label{exo:cyclo-poly}
Compute $\Phi_{12}(X)$ and verify that $\deg\Phi_{12}=\varphi(12)=4$.
\end{exercise}

\begin{exercise}[Normal closure]\label{exo:normal-closure}
Show that the normal closure of $\Q(\sqrt[3]{2}\,)/\Q$ in $\C$ is
$\Q(\sqrt[3]{2},\omega)$, and compute its Galois group.
\end{exercise}

\begin{exercise}[Non-Galois extension]\label{exo:non-galois}
Prove that $\Q(\sqrt[3]{2}\,)/\Q$ is not a Galois extension by showing
$|\!\Aut(\Q(\sqrt[3]{2}\,)/\Q)|=1<3=[\Q(\sqrt[3]{2}):\Q]$.
\end{exercise}

\begin{exercise}[Galois group over $\F_{p}$]\label{exo:gal-fp}
Let $f\in\F_{p}[X]$ be an irreducible polynomial of degree $n$.  Show that
the splitting field of $f$ is $\F_{p^{n}}$ and
$\Gal(\F_{p^{n}}/\F_{p})\cong\Z/n\Z$.
\end{exercise}

\begin{exercise}[Separability and Galois]\label{exo:sep-galois}
Give an example of a normal extension that is not Galois.
\textit{Hint:} work in characteristic $p$.
\end{exercise}

%% --- CHAPTER SUMMARY ---
\subsection*{Chapter Summary}
\addcontentsline{toc}{section}{Chapter Summary}

\begin{itemize}
\item A finite extension $L/K$ is \textbf{Galois} if and only if it is normal
and separable; equivalently, $L$ is the splitting field of a separable
polynomial, or $|\!\Aut(L/K)|=[L:K]$, or $L^{\Aut(L/K)}=K$.
\item \textbf{Artin's lemma}: for a finite $G\leq\Aut(L)$, $[L:L^{G}]=|G|$.
\item \textbf{Frobenius}: $\Gal(\F_{p^{n}}/\F_{p})=\langle x\mapsto
x^{p}\rangle\cong\Z/n\Z$.
\item \textbf{Cyclotomic}: $\Gal(\Q(\zeta_{n})/\Q)\cong(\Z/n\Z)^{\times}$.
\end{itemize}


%% ============================================================
%%  CHAPTER 12 — THE FUNDAMENTAL THEOREM OF GALOIS THEORY
%% ============================================================
\chapter{Galois Theory — The Fundamental Theorem}
\label{ch:fundamental-theorem}
\index{Fundamental Theorem of Galois Theory}

The Fundamental Theorem of Galois Theory establishes a precise dictionary
between the internal structure of a Galois group and the lattice of
intermediate fields.  It is one of the crowning achievements of algebra and
has profound consequences across mathematics.

%% ------------------------------------------------------------
\section{Statement and Proof of the Fundamental Theorem}
\label{sec:ftgt}

Let $L/K$ be a finite Galois extension with Galois group $G=\Gal(L/K)$.
Denote by
\[
  \mathcal{F} \;=\; \{E \mid K\subseteq E\subseteq L\}
\]
the set of intermediate fields, and by
\[
  \mathcal{G} \;=\; \{H \mid H\leq G\}
\]
the set of subgroups of $G$.

\begin{theorem}[Fundamental Theorem of Galois Theory]\label{thm:ftgt}
\index{Fundamental Theorem of Galois Theory}
Let $L/K$ be a finite Galois extension with Galois group $G=\Gal(L/K)$.
\begin{enumerate}
\item\label{it:ftgt-bij}
\textbf{(Galois correspondence).}\;
The maps
\begin{align*}
  \Phi\colon \mathcal{F}&\longrightarrow \mathcal{G},
    &E&\longmapsto \Gal(L/E),\\
  \Psi\colon \mathcal{G}&\longrightarrow \mathcal{F},
    &H&\longmapsto L^{H}=\Fix(H),
\end{align*}
are inclusion-reversing bijections, inverse to each other:
$\Psi\circ\Phi=\mathrm{id}_{\mathcal{F}}$ and
$\Phi\circ\Psi=\mathrm{id}_{\mathcal{G}}$.

\item\label{it:ftgt-deg}
\textbf{(Degree--index formula).}\;
For every intermediate field $E$ and the corresponding subgroup
$H=\Gal(L/E)$:
\[
  [L:E] \;=\; |H|, \qquad
  [E:K] \;=\; [G:H] \;=\; |G|/|H|.
\]

\item\label{it:ftgt-normal}
\textbf{(Normality criterion).}\;
An intermediate field $E$ is normal over $K$ (equivalently, $E/K$ is
Galois) if and only if the corresponding subgroup $H=\Gal(L/E)$ is a
normal subgroup of $G$.  In that case,
\[
  \Gal(E/K) \;\cong\; G/H.
\]
\end{enumerate}
\end{theorem}

\begin{proof}
We prove each part in turn.

\medskip\noindent
\textbf{Part~\ref{it:ftgt-bij}: The bijection.}

\smallskip\noindent
\emph{$\Psi\circ\Phi=\mathrm{id}$:}\;
Let $E\in\mathcal{F}$.  We must show $L^{\Gal(L/E)}=E$.  Since $L/K$ is
Galois (normal and separable), and $E$ is an intermediate field, $L/E$ is
also separable.  Moreover $L$ is a splitting field of some separable
polynomial over $K$, hence also over $E$, so $L/E$ is normal.  Therefore
$L/E$ is Galois, and the characterisation
(Theorem~\ref{thm:galois-equiv},
$\ref{it:gal-fix}$) gives $L^{\Gal(L/E)}=E$.

\smallskip\noindent
\emph{$\Phi\circ\Psi=\mathrm{id}$:}\;
Let $H\in\mathcal{G}$.  We must show $\Gal(L/L^{H})=H$.  By Artin's lemma
(Lemma~\ref{lem:artin}), $L/L^{H}$ is Galois with
$\Gal(L/L^{H})=H$.

\smallskip
Since both compositions are the identity, $\Phi$ and $\Psi$ are mutually
inverse bijections.  The inclusion-reversing property is immediate:
$E_{1}\subseteq E_{2}$ implies $\Gal(L/E_{2})\leq\Gal(L/E_{1})$, and
$H_{1}\leq H_{2}$ implies $L^{H_{2}}\subseteq L^{H_{1}}$.

\medskip\noindent
\textbf{Part~\ref{it:ftgt-deg}: Degree and index.}

\smallskip
Let $H=\Gal(L/E)$.  Since $L/E$ is Galois (as shown above),
Theorem~\ref{thm:gal-order} gives $[L:E]=|H|$.  By the tower law,
\[
  [E:K] \;=\; \frac{[L:K]}{[L:E]} \;=\; \frac{|G|}{|H|}
  \;=\; [G:H].
\]

\medskip\noindent
\textbf{Part~\ref{it:ftgt-normal}: Normal subgroups and normal extensions.}

\smallskip\noindent
$(\Longrightarrow)$\;
Suppose $E/K$ is normal (hence Galois, since $L/K$ separable implies $E/K$
separable).  We show $H=\Gal(L/E)\trianglelefteq G$.  Let $\sigma\in G$
and $\tau\in H$; we need $\sigma\tau\sigma^{-1}\in H$, i.e.,
$\sigma\tau\sigma^{-1}$ fixes $E$ pointwise.  For $a\in E$, since $E/K$ is
normal and $\sigma|_{E}$ is a $K$-embedding of $E$ into $L$, we have
$\sigma(E)=E$ (a normal extension is stable under any $K$-embedding).
So $\sigma^{-1}(a)\in E$, hence $\tau(\sigma^{-1}(a))=\sigma^{-1}(a)$
(because $\tau\in\Gal(L/E)$), and therefore
$\sigma\tau\sigma^{-1}(a)=\sigma(\sigma^{-1}(a))=a$.  Thus
$\sigma\tau\sigma^{-1}\in H$.

Moreover, the restriction map $G\to\Gal(E/K)$, $\sigma\mapsto\sigma|_{E}$,
is a well-defined surjective homomorphism (surjectivity uses the fact that
any $K$-automorphism of $E$ extends to $L$ since $L/E$ is Galois).  Its
kernel is $H$, so by the first isomorphism theorem,
$\Gal(E/K)\cong G/H$.

\smallskip\noindent
$(\Longleftarrow)$\;
Suppose $H\trianglelefteq G$.  We show $E=L^{H}$ is normal over $K$.
Let $\alpha\in E$ and $f=\irr(\alpha,K)$.  If $\beta$ is any root of $f$
in $L$, there exists $\sigma\in G$ with $\sigma(\alpha)=\beta$ (since $L/K$
is normal and $f$ is irreducible).  For any $\tau\in H$ we have
$\sigma^{-1}\tau\sigma\in H$ (since $H$ is normal), so
\[
  \tau(\beta)=\tau(\sigma(\alpha))
  =\sigma\bigl(\sigma^{-1}\tau\sigma(\alpha)\bigr)
  =\sigma(\alpha)=\beta,
\]
since $\sigma^{-1}\tau\sigma\in H$ fixes $\alpha\in E=L^{H}$.  Thus
$\beta\in L^{H}=E$, and $f$ splits in $E$.  Therefore $E/K$ is normal.
\end{proof}

%% --- SUMMARY DIAGRAM ---
\begin{center}
\begin{tikzcd}[row sep=2.5em, column sep=3em]
  \mathcal{F} \arrow[r,shift left=0.5ex,"\Phi"]
  & \mathcal{G} \arrow[l,shift left=0.5ex,"\Psi"]
\end{tikzcd}
\qquad
$\begin{aligned}
  E &\longmapsto \Gal(L/E) \\
  L^{H} &\longmapsfrom H
\end{aligned}$
\captionof{figure}{The Galois correspondence: intermediate fields $\leftrightarrow$
subgroups, inclusion-reversing.}\label{fig:galois-correspondence}
\end{center}

%% ------------------------------------------------------------
\section{Worked Examples with Complete Lattice Computations}
\label{sec:ftgt-examples}

\subsection{$\Q(\sqrt{2},\sqrt{3}\,)/\Q$: The Klein Four-Group}
\label{subsec:ex-v4}
\index{Galois correspondence!example with $V_4$}

\begin{example}[$\Q(\sqrt{2},\sqrt{3}\,)/\Q$ and $V_{4}$]
\label{ex:v4}
Let $L=\Q(\sqrt{2},\sqrt{3}\,)$.  Since $L$ is the splitting field of
$(X^{2}-2)(X^{2}-3)\in\Q[X]$ (separable), $L/\Q$ is Galois of degree $4$.
Define:
\begin{align*}
  \sigma &\colon \sqrt{2}\mapsto -\sqrt{2},\;\sqrt{3}\mapsto\sqrt{3}, \\
  \tau   &\colon \sqrt{2}\mapsto \sqrt{2},\;\sqrt{3}\mapsto -\sqrt{3}.
\end{align*}
Then $G=\Gal(L/\Q)=\{1,\sigma,\tau,\sigma\tau\}\cong V_{4}=(\Z/2\Z)^{2}$.

The subgroups and corresponding fixed fields are:

\[
\renewcommand{\arraystretch}{1.3}
\begin{array}{c|c|c}
  \text{Subgroup } H & \text{Fixed field } L^{H} & [L^{H}:\Q] \\\hline
  \{1\} & L=\Q(\sqrt{2},\sqrt{3}\,) & 4 \\
  \langle\sigma\rangle & \Q(\sqrt{3}\,) & 2 \\
  \langle\tau\rangle & \Q(\sqrt{2}\,) & 2 \\
  \langle\sigma\tau\rangle & \Q(\sqrt{6}\,) & 2 \\
  G & \Q & 1
\end{array}
\]

All subgroups of $V_{4}$ are normal (since $V_{4}$ is abelian), and
correspondingly every intermediate field is Galois over $\Q$.

\begin{center}
\begin{tikzcd}[row sep=1.8em, column sep=1em]
  & L=\Q(\sqrt{2},\sqrt{3}\,) \arrow[dl,dash] \arrow[d,dash]
    \arrow[dr,dash] & & &
  \{1\} & & \\
  \Q(\sqrt{2}\,) \arrow[dr,dash] & \Q(\sqrt{6}\,) \arrow[d,dash] &
  \Q(\sqrt{3}\,) \arrow[dl,dash] & &
  \langle\tau\rangle \arrow[u,dash] \arrow[dr,dash] &
  \langle\sigma\tau\rangle \arrow[ul,dash] \arrow[d,dash] &
  \langle\sigma\rangle \arrow[ull,dash] \arrow[dl,dash] \\
  & \Q & & & & G=V_{4} \arrow[from=1-5,dash,crossing over] &
\end{tikzcd}
\captionof{figure}{Lattice of intermediate fields (left) and
subgroups (right) for $\Q(\sqrt{2},\sqrt{3}\,)/\Q$.  The correspondence
reverses inclusions.}\label{fig:v4-lattice}
\end{center}
\end{example}

\subsection{$\Q(\sqrt[3]{2},\omega)/\Q$: The Symmetric Group $S_{3}$}
\label{subsec:ex-s3-lattice}
\index{Galois correspondence!example with $S_3$}

\begin{example}[{$\Q(\sqrt[3]{2},\omega)/\Q$ and $S_{3}$}]
\label{ex:s3-lattice}
From Example~\ref{ex:gal-s3}, $G=\Gal(\Q(\sqrt[3]{2},\omega)/\Q)\cong S_{3}$,
with $|G|=6$ and $[L:\Q]=6$.  Recall $\rho=(1\;2\;3)$ and $\tau=(2\;3)$
(in terms of the roots $\alpha_{1},\alpha_{2},\alpha_{3}$ of $X^{3}-2$).

The subgroups of $S_{3}$ and corresponding fixed fields:
\[
\renewcommand{\arraystretch}{1.3}
\begin{array}{c|c|c|c}
  \text{Subgroup } H & |H| & L^{H} & [L^{H}:\Q] \\\hline
  \{1\} & 1 & L & 6 \\
  \langle\tau\rangle=\{1,(2\;3)\} & 2 & \Q(\sqrt[3]{2}\,) & 3 \\
  \langle\rho\tau\rangle=\{1,(1\;3)\} & 2 & \Q(\omega\sqrt[3]{2}\,) & 3 \\
  \langle\rho^{2}\tau\rangle=\{1,(1\;2)\} & 2 & \Q(\omega^{2}\sqrt[3]{2}\,) & 3 \\
  \langle\rho\rangle=\{1,(1\;2\;3),(1\;3\;2)\} & 3 & \Q(\omega) & 2 \\
  G=S_{3} & 6 & \Q & 1
\end{array}
\]

The only normal subgroups of $S_{3}$ are $\{1\}$, $\langle\rho\rangle
\cong\Z/3\Z$, and $S_{3}$ itself.  Correspondingly, the only intermediate
fields that are Galois over $\Q$ are $L$, $\Q(\omega)$, and $\Q$.  Indeed:
\begin{itemize}
\item $\Q(\omega)/\Q$ is Galois (splitting field of $X^{2}+X+1$) with
      $\Gal(\Q(\omega)/\Q)\cong S_{3}/\langle\rho\rangle\cong\Z/2\Z$.
\item $\Q(\sqrt[3]{2}\,)/\Q$ is \emph{not} normal (hence not Galois),
      consistent with $\langle\tau\rangle$ not being normal in $S_{3}$.
\end{itemize}

\begin{center}
\begin{tikzcd}[row sep=1.8em, column sep=0.5em]
  & L \arrow[dl,dash,"2"'] \arrow[d,dash,"2"] \arrow[dr,dash,"2"]
    \arrow[drr,dash,"3"] & & \\
  \Q(\sqrt[3]{2}\,) \arrow[drr,dash,"3"'] &
  \Q(\omega\sqrt[3]{2}\,) \arrow[dr,dash,"3"'] &
  \Q(\omega^{2}\sqrt[3]{2}\,) \arrow[d,dash,"3"] &
  \Q(\omega) \arrow[dl,dash,"2"] \\
  & & \Q &
\end{tikzcd}
\qquad
\begin{tikzcd}[row sep=1.8em, column sep=0.5em]
  & & \{1\} \arrow[dll,dash] \arrow[dl,dash] \arrow[d,dash]
    \arrow[dr,dash] & \\
  \langle(2\;3)\rangle \arrow[dr,dash] &
  \langle(1\;3)\rangle \arrow[d,dash] &
  \langle(1\;2)\rangle \arrow[dl,dash] &
  \langle\rho\rangle \arrow[dll,dash] \\
  & S_{3} & &
\end{tikzcd}
\captionof{figure}{Lattice of intermediate fields (left) and subgroups (right) for
$\Q(\sqrt[3]{2},\omega)/\Q$.  Edge labels indicate degrees/indices.}
\label{fig:s3-lattice}
\end{center}
\end{example}

\subsection{$\F_{p^{12}}/\F_{p}$: Cyclic Galois Group}
\label{subsec:ex-fp12}
\index{Galois correspondence!example with finite fields}

\begin{example}[$\F_{p^{12}}/\F_{p}$ and $\Z/12\Z$]\label{ex:fp12}
By Theorem~\ref{thm:gal-finite-field},
$\Gal(\F_{p^{12}}/\F_{p})=\langle\varphi\rangle\cong\Z/12\Z$, where
$\varphi$ is the Frobenius.

The subgroups of $\Z/12\Z$ are $\langle d\rangle\cong\Z/(12/d)\Z$ for each
divisor $d$ of $12$: $d\in\{1,2,3,4,6,12\}$.  The fixed field of
$\langle\varphi^{d}\rangle$ is $\F_{p^{d}}$.

\[
\renewcommand{\arraystretch}{1.3}
\begin{array}{c|c|c}
  \text{Subgroup } \langle\varphi^{d}\rangle & \text{Order} &
  \text{Fixed field} \\\hline
  \langle\varphi\rangle & 12 & \F_{p} \\
  \langle\varphi^{2}\rangle & 6 & \F_{p^{2}} \\
  \langle\varphi^{3}\rangle & 4 & \F_{p^{3}} \\
  \langle\varphi^{4}\rangle & 3 & \F_{p^{4}} \\
  \langle\varphi^{6}\rangle & 2 & \F_{p^{6}} \\
  \{1\} & 1 & \F_{p^{12}}
\end{array}
\]

Since $\Z/12\Z$ is abelian, every subgroup is normal, so every
$\F_{p^{d}}/\F_{p}$ is Galois — as expected.

\begin{center}
\begin{tikzcd}[row sep=1.5em, column sep=0.8em]
  & \F_{p^{12}} \arrow[dl,dash,"2"'] \arrow[d,dash,"3"]
    \arrow[dr,dash,"2"] & \\
  \F_{p^{6}} \arrow[d,dash,"2"'] \arrow[dr,dash,"3"] &
  \F_{p^{4}} \arrow[dl,dash,"2"] \arrow[dr,dash,"2"] &
  \F_{p^{3}} \arrow[dl,dash] \arrow[d,dash,"3"] \\
  \F_{p^{3}} \arrow[dr,dash,"3"'] & \F_{p^{2}} \arrow[d,dash,"2"] &
  \F_{p} \\
  & \F_{p} &
\end{tikzcd}
\captionof{figure}{Subfield lattice of $\F_{p^{12}}/\F_{p}$.  The divisibility
lattice of $12$ governs the structure.}\label{fig:fp12-lattice}
\end{center}
\end{example}

%% ------------------------------------------------------------
\section{The Primitive Element Theorem}
\label{sec:primitive-element}
\index{Primitive Element Theorem}

\begin{theorem}[Primitive Element Theorem]\label{thm:primitive-element}
Let $L/K$ be a finite separable extension.  Then there exists $\alpha\in L$
such that $L=K(\alpha)$.  Such an $\alpha$ is called a \emph{primitive
element}\index{primitive element} of the extension.
\end{theorem}

\begin{proof}
If $K$ is finite, then $L$ is also finite, and $L^{\times}$ is cyclic
(a finite subgroup of the multiplicative group of a field is cyclic).  Any
generator $\alpha$ of $L^{\times}$ satisfies $L=K(\alpha)$.

Assume $K$ is infinite.  It suffices to treat the case $L=K(\beta,\gamma)$
(the general case follows by induction).  Let
$f=\irr(\beta,K)$ with roots $\beta=\beta_{1},\dots,\beta_{m}$ and
$g=\irr(\gamma,K)$ with roots $\gamma=\gamma_{1},\dots,\gamma_{n}$ in
$\overline{K}$.  Since $L/K$ is separable, the roots of $f$ and $g$ are
distinct.

For $i\geq 1$ and $j\geq 2$, the equation
$\beta_{i}+t\gamma_{j}=\beta+t\gamma$ (equivalently
$t=(\beta_{i}-\beta)/(\gamma-\gamma_{j})$) has at most one solution
$t\in K$.  Since $K$ is infinite, we may choose $c\in K$ avoiding all such
values.  Set $\alpha=\beta+c\gamma$.  We claim $L=K(\alpha)$.

Clearly $K(\alpha)\subseteq K(\beta,\gamma)=L$.  For the reverse, set
$h(X)=f(\alpha-cX)\in K(\alpha)[X]$.  Then $h(\gamma)=f(\beta)=0$, so
$\gamma$ is a common root of $h(X)$ and $g(X)$ in $\overline{K}$.  We show
$\gcd(h,g)$ in $K(\alpha)[X]$ has $\gamma$ as its only root.

If $\gamma_{j}$ (with $j\geq 2$) were also a root of $h$, then
$f(\alpha-c\gamma_{j})=0$, so $\alpha-c\gamma_{j}=\beta_{i}$ for some $i$.
Then $\beta+c\gamma-c\gamma_{j}=\beta_{i}$, i.e.,
$c=(\beta_{i}-\beta)/(\gamma-\gamma_{j})$, contradicting our choice of $c$.

Therefore $\gcd(h,g)$ in $K(\alpha)[X]$ is the minimal polynomial of
$\gamma$ over $K(\alpha)$, which must be linear, giving
$\gamma\in K(\alpha)$.  Hence $\beta=\alpha-c\gamma\in K(\alpha)$, and
$L=K(\beta,\gamma)\subseteq K(\alpha)$.
\end{proof}

\begin{example}[Primitive element of $\Q(\sqrt{2},\sqrt{3}\,)$]
\label{ex:primitive-element}
We have $\Q(\sqrt{2},\sqrt{3}\,)=\Q(\sqrt{2}+\sqrt{3}\,)$.  Indeed,
if $\alpha=\sqrt{2}+\sqrt{3}$, then $\alpha^{2}=5+2\sqrt{6}$, so
$\sqrt{6}\in\Q(\alpha)$, whence $\sqrt{3}=\alpha\sqrt{6}-3\alpha=
(\alpha^{2}-5)/2\cdot\alpha^{-1}$... more directly: $\alpha^{2}=5+2\sqrt{6}$
gives $\sqrt{6}\in\Q(\alpha)$; then $\sqrt{3}=\sqrt{6}/\sqrt{2}$ and
$\sqrt{2}=\alpha-\sqrt{3}\in\Q(\alpha)$.
\end{example}

%% ------------------------------------------------------------
\section{Exercises}
\label{sec:ch12-exercises}

\begin{exercise}[Full Galois correspondence for $\Q(\sqrt{2},\sqrt{3},\sqrt{5}\,)/\Q$]
\label{exo:three-sqrt}
Show that $\Gal(\Q(\sqrt{2},\sqrt{3},\sqrt{5}\,)/\Q)\cong(\Z/2\Z)^{3}$
and list all $16$ subgroups with their fixed fields.
\end{exercise}

\begin{exercise}[Galois correspondence for $X^{4}-2$]\label{exo:x4-2-lattice}
Draw the complete lattice of subgroups of $D_{4}$ and the corresponding
lattice of intermediate fields for the splitting field of $X^{4}-2$ over
$\Q$.  Identify which intermediate extensions are Galois over $\Q$.
\end{exercise}

\begin{exercise}[The FTGT for cyclic extensions]\label{exo:ftgt-cyclic}
Let $L/K$ be a Galois extension with $\Gal(L/K)\cong\Z/n\Z$.  Prove that
the intermediate fields are in bijection with the divisors of $n$, and that
every intermediate extension is Galois over $K$.
\end{exercise}

\begin{exercise}[Degree formula verification]\label{exo:deg-formula}
In Example~\ref{ex:s3-lattice}, verify the degree--index formula
$[E:K]=[G:H]$ for each intermediate field $E$.
\end{exercise}

\begin{exercise}[Non-abelian Galois group]\label{exo:non-abelian}
Give an example of a Galois extension $L/K$ with an intermediate field $E$
such that $E/K$ is \emph{not} Galois.
\end{exercise}

\begin{exercise}[Galois closure]\label{exo:galois-closure}
Let $L/K$ be a finite separable extension.  Prove that there exists a
smallest Galois extension $M/K$ containing $L$, and express $\Gal(M/K)$
as a subgroup of a symmetric group.
\end{exercise}

\begin{exercise}[Intersection and compositum]\label{exo:intersection-compositum}
Let $L/K$ be Galois with $G=\Gal(L/K)$, and let $E_{1},E_{2}$ be
intermediate fields with corresponding subgroups $H_{1},H_{2}$.  Show:
\begin{enumerate}
\item $E_{1}\cap E_{2}\longleftrightarrow\langle H_{1},H_{2}\rangle$.
\item $E_{1}E_{2}\longleftrightarrow H_{1}\cap H_{2}$.
\end{enumerate}
\end{exercise}

\begin{exercise}[{Primitive element for $\Q(\sqrt[4]{2},i)$}]\label{exo:primitive-x4}
Find a primitive element for $\Q(\sqrt[4]{2},i)/\Q$.
\end{exercise}

\begin{exercise}[Finite fields and the FTGT]\label{exo:ftgt-finite}
Apply the Fundamental Theorem to $\F_{p^{30}}/\F_{p}$.  List all
intermediate fields and draw the lattice.
\end{exercise}

\begin{exercise}[Abelian extensions of $\Q$]\label{exo:abelian-ext}
Show that every finite abelian extension of $\Q$ is contained in some
cyclotomic extension $\Q(\zeta_{n})$.
\textit{(This is the Kronecker--Weber theorem; give a proof sketch.)}
\end{exercise}

\begin{exercise}[Quotient groups and intermediate Galois groups]
\label{exo:quotient-gal}
Let $L/K$ be Galois and $E$ an intermediate field with $E/K$ Galois.
Show that the restriction map $\Gal(L/K)\to\Gal(E/K)$ is surjective with
kernel $\Gal(L/E)$, recovering $\Gal(E/K)\cong\Gal(L/K)/\Gal(L/E)$.
\end{exercise}

%% --- CHAPTER SUMMARY ---
\subsection*{Chapter Summary}
\addcontentsline{toc}{section}{Chapter Summary}

\begin{itemize}
\item The \textbf{Fundamental Theorem of Galois Theory} establishes an
inclusion-reversing bijection between intermediate fields of a Galois
extension $L/K$ and subgroups of $\Gal(L/K)$.
\item \textbf{Degree--index duality}: $[E:K]=[G:\Gal(L/E)]$.
\item \textbf{Normality criterion}: $E/K$ is Galois $\Leftrightarrow$
$\Gal(L/E)\trianglelefteq\Gal(L/K)$, and then
$\Gal(E/K)\cong\Gal(L/K)/\Gal(L/E)$.
\item The \textbf{Primitive Element Theorem} guarantees that every finite
separable extension is simple: $L=K(\alpha)$.
\end{itemize}


%% ============================================================
%%  CHAPTER 13 — APPLICATIONS
%% ============================================================
\chapter{Applications of Galois Theory}\label{ch:applications}

In this final chapter we present three celebrated applications of Galois
theory: the classical Greek construction problems, the solvability of
polynomial equations by radicals, and a proof of the Fundamental Theorem of
Algebra.

%% ============================================================
\section{Ruler-and-Compass Constructibility}
\label{sec:constructibility}
\index{ruler-and-compass construction}\index{constructible number}

\subsection{Constructible Numbers}
\label{subsec:constructible-numbers}

Starting from two marked points at distance $1$, a real number $\alpha$ is
\emph{constructible} if a segment of length $|\alpha|$ can be produced in
finitely many steps using an unmarked straightedge and a compass.

\begin{definition}[Field of constructible numbers]\label{def:constructible-field}
\index{constructible number}
A real number $\alpha$ is \emph{constructible} if it can be obtained from
$\{0,1\}$ by a finite sequence of additions, subtractions, multiplications,
divisions, and extractions of square roots of positive numbers.
\end{definition}

\begin{theorem}[Constructible numbers form a field]\label{thm:constructible-field}
The set of constructible numbers forms a subfield of $\R$, closed under
taking square roots of positive elements.
\end{theorem}

\begin{proof}
Each basic straightedge-and-compass operation corresponds to:
\begin{itemize}
\item intersecting two lines (solving a linear system — yields elements of
      the current field);
\item intersecting a line with a circle, or two circles (solving a quadratic
      equation — yields elements in a degree~$\leq 2$ extension of the
      current field).
\end{itemize}
Hence the set of constructible numbers is closed under $+,-,\times,\div$
and $\sqrt{\phantom{x}}$ (for positive elements).  Closure under the
four arithmetic operations shows it is a field.
\end{proof}

\begin{theorem}[Degree criterion for constructibility]\label{thm:degree-constructible}
\index{constructible number!degree criterion}
If $\alpha\in\R$ is constructible, then $[\Q(\alpha):\Q]=2^{s}$ for some
$s\geq 0$.
\end{theorem}

\begin{proof}
Starting from $K_{0}=\Q$, each construction step produces a point in a
field $K_{i+1}$ with $[K_{i+1}:K_{i}]\in\{1,2\}$ (line--line
intersections stay in the same field; line--circle or circle--circle
intersections adjoin at most a square root).  After $r$ steps,
$\alpha\in K_{r}$ where
\[
  [K_{r}:\Q] = [K_{r}:K_{r-1}]\cdots[K_{1}:K_{0}] = 2^{s}
\]
for some $s\leq r$.  By the tower law, $[\Q(\alpha):\Q]$ divides
$[K_{r}:\Q]=2^{s}$, so $[\Q(\alpha):\Q]$ is also a power of $2$.
\end{proof}

\begin{remark}[Converse]\label{rem:constructible-converse}
The converse is not quite true in general: $[\Q(\alpha):\Q]$ being a power
of $2$ is necessary but not sufficient.  However, $\alpha$ is constructible
if and only if $[\Q(\alpha):\Q]$ is a power of $2$ \emph{and} the Galois
closure of $\Q(\alpha)/\Q$ has degree a power of $2$ over $\Q$ (equivalently,
the Galois group is a $2$-group).
\end{remark}

\subsection{The Three Classical Impossibilities}
\label{subsec:impossibilities}
\index{doubling the cube}\index{trisecting the angle}\index{squaring the circle}

\begin{corollary}[Doubling the cube is impossible]\label{cor:doubling-cube}
It is impossible to construct, with ruler and compass alone, the side of a
cube of volume~$2$.
\end{corollary}

\begin{proof}
The required length is $\sqrt[3]{2}$.  We have $\irr(\sqrt[3]{2},\Q)=X^{3}-2$
(irreducible by Eisenstein at $p=2$), so $[\Q(\sqrt[3]{2}):\Q]=3$.  Since
$3$ is not a power of $2$, the construction is impossible by
Theorem~\ref{thm:degree-constructible}.
\end{proof}

\begin{corollary}[Trisecting a general angle is impossible]\label{cor:trisecting}
There is no ruler-and-compass construction that trisects every given angle.
\end{corollary}

\begin{proof}
It suffices to show that the angle $60°$ cannot be trisected, i.e., that
$\cos 20°$ is not constructible.  The triple-angle formula
$\cos 3\theta=4\cos^{3}\theta-3\cos\theta$ with $\theta=20°$ gives
$4c^{3}-3c=\cos 60°=1/2$, where $c=\cos 20°$.  Setting $c=t/2$:
$t^{3}-3t-1=0$.  This polynomial is irreducible over $\Q$ (it has no
rational roots by the rational root theorem), so $[\Q(\cos 20°):\Q]=3$.
Again $3\neq 2^{s}$, so the construction is impossible.
\end{proof}

\begin{corollary}[Squaring the circle is impossible]\label{cor:squaring-circle}
It is impossible to construct with ruler and compass a square of area $\pi$.
\end{corollary}

\begin{proof}
The required side length is $\sqrt{\pi}$, so it suffices to show $\pi$ is
not constructible.  But $\pi$ is transcendental (Lindemann, 1882), hence
$[\Q(\pi):\Q]=\infty$, which is certainly not a power of $2$.
\end{proof}

\subsection{Constructibility of Regular Polygons}
\label{subsec:regular-polygons}
\index{regular polygon!constructibility}\index{Gauss--Wantzel theorem}
\index{Fermat prime}

\begin{theorem}[Gauss--Wantzel]\label{thm:gauss-wantzel}
A regular $n$-gon (with $n\geq 3$) is constructible with ruler and compass
if and only if
\[
  n = 2^{a}\,p_{1}\,p_{2}\cdots p_{r},
\]
where $a\geq 0$ and $p_{1},\dots,p_{r}$ are \emph{distinct} Fermat primes,
i.e., primes of the form $F_{k}=2^{2^{k}}+1$.
\end{theorem}

\begin{proof}[Proof sketch]
Constructing a regular $n$-gon amounts to constructing $\cos(2\pi/n)$.
One has $\Q(\cos(2\pi/n))\subseteq\Q(\zeta_{n})$ with
$[\Q(\zeta_{n}):\Q(\cos(2\pi/n))]=2$.  By
Theorem~\ref{thm:gal-cyclotomic},
$[\Q(\zeta_{n}):\Q]=\varphi(n)$.  The angle $2\pi/n$ is constructible if
and only if $\varphi(n)$ is a power of $2$, which holds precisely when $n$
has the given form.
\end{proof}

\begin{remark}[Known Fermat primes]\label{rem:fermat-primes}
The known Fermat primes are $F_{0}=3$, $F_{1}=5$, $F_{2}=17$, $F_{3}=257$,
$F_{4}=65537$.  It is unknown whether any others exist.  Thus the
constructible regular $n$-gons for odd $n$ include $n=3,5,15,17,51,85,
255,257,\dots$
\end{remark}

%% --- TikZ construction illustration ---
\begin{center}
\begin{tikzpicture}[scale=1.8]
  % Regular pentagon
  \foreach \i in {0,...,4} {
    \coordinate (P\i) at ({90+72*\i}:1);
  }
  \draw[thick] (P0)--(P1)--(P2)--(P3)--(P4)--cycle;
  \foreach \i in {0,...,4} {
    \fill (P\i) circle (1pt);
  }
  \node at (0,-1.4) {Regular pentagon ($n=5=F_1$)};
\end{tikzpicture}
\hspace{2cm}
\begin{tikzpicture}[scale=1.8]
  % Regular heptagon with "impossible" label
  \foreach \i in {0,...,6} {
    \coordinate (Q\i) at ({90+360/7*\i}:1);
  }
  \draw[thick,dashed,red!70!black] (Q0)--(Q1)--(Q2)--(Q3)--(Q4)--(Q5)--(Q6)--cycle;
  \foreach \i in {0,...,6} {
    \fill[red!70!black] (Q\i) circle (1pt);
  }
  \node[red!70!black] at (0,-1.4) {Regular heptagon ($n=7$): impossible};
\end{tikzpicture}
\captionof{figure}{The regular pentagon is constructible ($5$ is a Fermat prime);
the regular heptagon is not ($\varphi(7)=6\neq 2^{s}$).}
\label{fig:polygons}
\end{center}

%% ============================================================
\section{Solvability by Radicals}
\label{sec:solvability-radicals}
\index{solvability by radicals}\index{radical extension}

\subsection{Radical Extensions and Solvable Groups}
\label{subsec:radical-ext}

\begin{definition}[Radical extension]\label{def:radical-ext}
\index{radical extension}
A field extension $L/K$ is a \emph{radical extension} if there exists a
tower
\[
  K = K_{0} \;\subset\; K_{1} \;\subset\; \cdots \;\subset\; K_{r} = L
\]
where each $K_{i+1}=K_{i}(\alpha_{i})$ with $\alpha_{i}^{n_{i}}\in K_{i}$
for some $n_{i}\geq 1$.
\end{definition}

\begin{definition}[Solvable by radicals]\label{def:solvable-radicals}
\index{solvable by radicals}
A polynomial $f\in K[X]$ is \emph{solvable by radicals} if its splitting
field is contained in some radical extension of $K$.
\end{definition}

\begin{definition}[Solvable group]\label{def:solvable-group}
\index{solvable group}
A group $G$ is \emph{solvable} if there exists a chain of subgroups
\[
  \{1\} = G_{0} \;\trianglelefteq\; G_{1} \;\trianglelefteq\;
  \cdots \;\trianglelefteq\; G_{r} = G
\]
such that each quotient $G_{i+1}/G_{i}$ is abelian.
\end{definition}

\begin{example}[Solvable groups]\label{ex:solvable-groups}
\begin{enumerate}
\item Every abelian group is solvable (take $r=1$).
\item $S_{3}$ is solvable: $\{1\}\trianglelefteq A_{3}\trianglelefteq S_{3}$
      with $A_{3}/\{1\}\cong\Z/3\Z$ and $S_{3}/A_{3}\cong\Z/2\Z$.
\item $S_{4}$ is solvable:
      $\{1\}\trianglelefteq V_{4}\trianglelefteq A_{4}\trianglelefteq S_{4}$.
\end{enumerate}
\end{example}

\subsection{$S_{n}$ Is Not Solvable for $n\geq 5$}
\label{subsec:sn-not-solvable}

\begin{theorem}[$S_{n}$ is not solvable for $n\geq 5$]\label{thm:sn-not-solvable}
\index{symmetric group!not solvable}
For $n\geq 5$, the symmetric group $S_{n}$ is not solvable.
\end{theorem}

\begin{proof}
We first establish that $A_{n}$ is simple for $n\geq 5$, then use this
to prove $S_{n}$ is not solvable.

\medskip\noindent
\textbf{Step 1: $A_{n}$ is simple for $n\geq 5$.}\;
Let $N\trianglelefteq A_{n}$ with $N\neq\{1\}$.  We show $N=A_{n}$ by
proving $N$ contains all $3$-cycles (which generate $A_{n}$).

Let $\sigma\in N\setminus\{1\}$.  We consider cases based on the cycle
structure of $\sigma$.

\emph{Case 1:} $\sigma$ contains a cycle of length $\geq 3$, say
$\sigma$ moves $1\mapsto 2\mapsto 3$.  Let $\tau=(3\;4\;5)\in A_{n}$.
Then the commutator $[\sigma,\tau]=\sigma\tau\sigma^{-1}\tau^{-1}\in N$
is a non-identity element that can be shown (by direct computation,
considering sub-cases) to yield a permutation moving fewer elements than
$\sigma$.  Iterating, we reduce to elements that are products of fewer
cycles.

\emph{Case 2:} $\sigma$ is a product of disjoint transpositions, say
$\sigma=(1\;2)(3\;4)\cdots$.  Then for $\tau=(1\;2\;3)\in A_{n}$,
$\tau\sigma\tau^{-1}=(2\;3)(1\;4)\cdots\in N$, and multiplying:
$\sigma\cdot\tau\sigma\tau^{-1}=(1\;3\;4\;2)\cdots\in N$ (or a shorter
permutation), reducing to Case~1 or giving a $3$-cycle directly.

In all cases, since $n\geq 5$ provides enough ``room'' to manoeuvre with
conjugation, one shows that $N$ must contain a $3$-cycle.  Since all
$3$-cycles are conjugate in $A_{n}$ (for $n\geq 5$), $N$ contains all
$3$-cycles, hence $N=A_{n}$.

\medskip\noindent
\textbf{Step 2: $S_{n}$ is not solvable.}\;
Suppose $S_{n}$ is solvable with composition series
$\{1\}=G_{0}\trianglelefteq\cdots\trianglelefteq G_{r}=S_{n}$ (abelian
quotients).  Then $A_{n}\cap G_{i}$ gives a composition series for $A_{n}$
with abelian quotients, so $A_{n}$ would be solvable.  But the only normal
subgroups of $A_{n}$ (for $n\geq 5$) are $\{1\}$ and $A_{n}$, and
$A_{n}/\{1\}=A_{n}$ is not abelian for $n\geq 5$.  Contradiction.
\end{proof}

\subsection{Galois's Theorem and the Abel--Ruffini Theorem}
\label{subsec:galois-theorem}

\begin{theorem}[Galois's theorem]\label{thm:galois-solvability}
\index{Galois's theorem (solvability)}
Let $f\in K[X]$ be a separable polynomial with splitting field $L$ over $K$,
where $\chr(K)=0$.  Then $f$ is solvable by radicals if and only if
$\Gal(L/K)$ is a solvable group.
\end{theorem}

\begin{proof}[Proof sketch]
$(\Longrightarrow)$\;
If $f$ is solvable by radicals, the splitting field $L$ is contained in a
radical tower $K=K_{0}\subset K_{1}\subset\cdots\subset K_{r}$.  After
adjoining appropriate roots of unity, each step $K_{i}\subset K_{i+1}$
becomes a cyclic (hence abelian) Galois extension (Kummer theory).  The
Galois group $\Gal(L/K)$ is then a quotient of a group with an abelian
composition series, hence is solvable.

$(\Longleftarrow)$\;
If $\Gal(L/K)$ is solvable, one constructs a radical tower containing $L$
by reversing the argument: the composition series with abelian (in fact
cyclic) quotients corresponds, via Kummer theory and the Fundamental Theorem,
to successive radical adjunctions.
\end{proof}

\begin{theorem}[Abel--Ruffini]\label{thm:abel-ruffini}
\index{Abel--Ruffini theorem}
The general polynomial equation of degree $n\geq 5$ is not solvable by
radicals.  More precisely, for each $n\geq 5$ there exist polynomials
$f\in\Q[X]$ of degree $n$ whose Galois group over $\Q$ is $S_{n}$, and
since $S_{n}$ is not solvable (Theorem~\ref{thm:sn-not-solvable}), $f$ is
not solvable by radicals.
\end{theorem}

\begin{example}[$X^{5}-4X+2$ has Galois group $S_{5}$]\label{ex:x5-s5}
\index{Galois group!of $X^5-4X+2$}
Consider $f=X^{5}-4X+2\in\Q[X]$.

\begin{enumerate}
\item \textbf{Irreducibility.}\; By Eisenstein's criterion at $p=2$, $f$ is
irreducible over $\Q$.

\item \textbf{Real roots.}\; One checks that $f$ has exactly three real
roots and two complex conjugate roots (by analysing $f'=5X^{4}-4$: the
critical points are at $x=\pm(4/5)^{1/4}$; evaluating $f$ at these points
and at $\pm\infty$ shows three real zeros).

\item \textbf{Galois group.}\; The Galois group $G=\Gal(L/\Q)$ embeds into
$S_{5}$ (acting on the five roots).  Since $f$ is irreducible of degree $5$,
$G$ acts transitively, so $5\mid|G|$.  Complex conjugation gives a
transposition in $G$ (swapping the two complex roots).  A transitive
subgroup of $S_{5}$ containing a transposition and an element of order $5$
(by Cauchy's theorem, since $5\mid|G|$) is all of $S_{5}$.  Hence
$G\cong S_{5}$.
\end{enumerate}

Since $S_{5}$ is not solvable, $f$ is not solvable by radicals.
\end{example}

%% ============================================================
\section{The Fundamental Theorem of Algebra}
\label{sec:fta}
\index{Fundamental Theorem of Algebra}

\begin{theorem}[Fundamental Theorem of Algebra]\label{thm:fta}
The field $\C$ is algebraically closed: every non-constant polynomial
$f\in\C[X]$ has a root in $\C$.
\end{theorem}

\begin{proof}[Proof via Galois theory]
We use two analytic facts: (i)~every polynomial in $\R[X]$ of odd degree has
a real root; (ii)~every positive real number has a real square root.

Let $f\in\R[X]$ be irreducible.  Let $L$ be the splitting field of $f$ over
$\R$.  Then $L/\R$ is a finite Galois extension; set $G=\Gal(L/\R)$
and $|G|=2^{s}\cdot m$ with $m$ odd.

\medskip\noindent
\textbf{Step 1:} $m=1$ (i.e., $|G|$ is a power of $2$).

Let $P$ be a Sylow $2$-subgroup of $G$ and $E=L^{P}$.  Then
$[E:\R]=[G:P]=m$.  By the Primitive Element Theorem, $E=\R(\alpha)$ and
$\irr(\alpha,\R)$ has odd degree $m$.  By fact~(i), this polynomial has a
root in $\R$, so $m=1$ (since it is irreducible and has a real root, it
must be linear).  Hence $|G|=2^{s}$.

\medskip\noindent
\textbf{Step 2:} $s\leq 1$ (i.e., $|G|\leq 2$).

Suppose $s\geq 2$.  Since $G$ is a $2$-group, it has a subgroup $H$ of
index $2$ (indeed, of index $2$ in some subgroup of index $2$, etc.).
Let $M=L^{H}$.  Then $[M:\R]=2$, so $M=\R(\beta)$ with
$\irr(\beta,\R)=X^{2}+bX+c$ for some $b,c\in\R$.  By fact~(ii), the
discriminant $b^{2}-4c$ satisfies: if $b^{2}-4c<0$, then
$\beta=(-b\pm i\sqrt{4c-b^{2}})/2\in\C$.  But we are working with
$L/\R$, and $\C=\R(i)$ already contains all roots of quadratics over $\R$.
So $M\subseteq\C$; but $[M:\R]=2=[\C:\R]$, hence $M=\C$.

Now consider $\Gal(L/\C)$.  If $\Gal(L/\C)$ had a subgroup of index $2$,
the same argument would produce a degree-$2$ extension of $\C$, say
$\C(\gamma)$ with $\gamma^{2}+b'\gamma+c'=0$ for $b',c'\in\C$.  But the
quadratic formula shows $\gamma\in\C$ (since $\C$ is closed under square
roots by fact~(ii) and the polar form of complex numbers).  This is a
contradiction: the extension would be trivial.

Since $\Gal(L/\C)$ is a $2$-group that has no subgroup of index $2$ (other
than itself if it is non-trivial), we must have $\Gal(L/\C)=\{1\}$, i.e.,
$L=\C$.  Hence $|G|=[\C:\R]=2$ and $s=1$.

\medskip\noindent
\textbf{Conclusion.}\;
Every irreducible $f\in\R[X]$ has $\deg f\leq 2$.  Equivalently, every
$f\in\R[X]$ factors into linear and quadratic factors over $\R$, hence
into linear factors over $\C$.  It follows that $\C$ is algebraically
closed.
\end{proof}

%% ============================================================
\section{Concluding Remarks: Further Directions}
\label{sec:further}

Galois theory, far from being a closed chapter, serves as a gateway to
some of the deepest areas of modern mathematics.

\begin{itemize}
\item \textbf{Algebraic geometry.}\;
The Galois group of the function field of an algebraic curve encodes
arithmetic and geometric information.  \'Etale cohomology, developed by
Grothendieck, generalises Galois theory to schemes and is the foundation of
modern arithmetic geometry\index{algebraic geometry}.

\item \textbf{Algebraic number theory.}\;
The absolute Galois group $\Gal(\overline{\Q}/\Q)$ is the central object of
study.  Class field theory classifies abelian extensions of number fields via
Artin reciprocity, while the Langlands programme seeks a vast non-abelian
generalisation\index{number theory}.

\item \textbf{Inverse Galois problem.}\;
Which finite groups arise as $\Gal(L/\Q)$?  This remains one of the great
open questions.  Shafarevich proved every solvable group occurs; the general
case is open\index{inverse Galois problem}.

\item \textbf{Differential Galois theory.}\;
Replacing polynomial equations by linear differential equations, the
Picard--Vessiot theory develops a Galois correspondence for differential
fields, with applications to the integrability of dynamical
systems\index{differential Galois theory}.

\item \textbf{Representation theory.}\;
Representations of Galois groups (Artin representations, $\ell$-adic
representations) are at the heart of the Langlands
programme\index{representation theory}.
\end{itemize}

%% ============================================================
\section{Exercises}
\label{sec:ch13-exercises}

\begin{exercise}[Constructibility of $\cos(2\pi/17)$]\label{exo:17-gon}
Use the fact that $\Gal(\Q(\zeta_{17})/\Q)\cong(\Z/17\Z)^{\times}\cong
\Z/16\Z$ to explain why the regular $17$-gon is constructible.  Describe
the tower of quadratic extensions.
\end{exercise}

\begin{exercise}[{Non-constructibility of $\sqrt[3]{2}$}]\label{exo:cube-root-2}
Give a self-contained proof that $\sqrt[3]{2}$ is not constructible.
\end{exercise}

\begin{exercise}[Solvable Galois groups of cubics and quartics]
\label{exo:solvable-low}
Show that the Galois group of any irreducible polynomial of degree $\leq 4$
over $\Q$ is solvable, and conclude that every such polynomial is solvable
by radicals.
\end{exercise}

\begin{exercise}[Explicit radical solution of a cubic]\label{exo:cubic-radical}
Solve $X^{3}-3X+1=0$ by radicals.  \textit{Hint:} the discriminant is $81$,
a perfect square, so $\Gal\cong\Z/3\Z$.
\end{exercise}

\begin{exercise}[Galois group of $X^{5}-2$]\label{exo:x5-2}
Show that $\Gal(\Q(\sqrt[5]{2},\zeta_{5})/\Q)\cong\Z/5\Z\rtimes
(\Z/5\Z)^{\times}$ (the Frobenius group of order $20$), and conclude that
$X^{5}-2$ is solvable by radicals.
\end{exercise}

\begin{exercise}[Another unsolvable quintic]\label{exo:x5-6x-3}
Show that $X^{5}-6X+3\in\Q[X]$ has Galois group $S_{5}$ and is therefore
not solvable by radicals.
\end{exercise}

\begin{exercise}[FTA via Liouville]\label{exo:fta-liouville}
Give an alternative proof of the Fundamental Theorem of Algebra using
Liouville's theorem from complex analysis.
\end{exercise}

\begin{exercise}[Constructible regular polygons]\label{exo:constructible-ngon}
For which $n\leq 30$ is the regular $n$-gon constructible?  List all such $n$.
\end{exercise}

\begin{exercise}[Composition series of $S_4$]\label{exo:comp-s4}
Write down a composition series for $S_{4}$ and verify that all composition
factors are cyclic of prime order.
\end{exercise}

\begin{exercise}[Solvability and the discriminant]\label{exo:discriminant-solvable}
Let $f\in\Q[X]$ be irreducible of degree $5$ with discriminant $\Delta$.
Show that if $\sqrt{\Delta}\in\Q$, then $\Gal(f)\leq A_{5}$.  Since $A_{5}$
is not solvable, such a polynomial is still generally not solvable by
radicals.
\end{exercise}

%% --- CHAPTER SUMMARY ---
\subsection*{Chapter Summary}
\addcontentsline{toc}{section}{Chapter Summary}

\begin{itemize}
\item A real number $\alpha$ is \textbf{constructible} only if
$[\Q(\alpha):\Q]$ is a power of $2$.  This rules out doubling the cube,
trisecting a general angle, and squaring the circle.
\item \textbf{Gauss--Wantzel}: a regular $n$-gon is constructible
$\Leftrightarrow$ $n=2^{a}p_{1}\cdots p_{r}$ with distinct Fermat primes
$p_{i}$.
\item \textbf{Galois's theorem}: a polynomial (in characteristic $0$) is
solvable by radicals $\Leftrightarrow$ its Galois group is solvable.
\item Since $S_{n}$ is not solvable for $n\geq 5$, the \textbf{general
quintic} is not solvable by radicals (\textbf{Abel--Ruffini}).
\item The \textbf{Fundamental Theorem of Algebra} admits an elegant proof
using Galois theory and Sylow theory.
\end{itemize}


%% ============================================================
%%  BIBLIOGRAPHY
%% ============================================================
\begin{thebibliography}{99}

\bibitem{dummit-foote}
D.~S. Dummit and R.~M. Foote,
\textit{Abstract Algebra}, 3rd ed.,
John Wiley \& Sons, Hoboken, NJ, 2004.
\index{Dummit--Foote}

\bibitem{lang}
S.~Lang,
\textit{Algebra}, rev.\ 3rd ed.,
Graduate Texts in Mathematics 211, Springer-Verlag, New York, 2002.
\index{Lang}

\bibitem{artin}
M.~Artin,
\textit{Algebra}, 2nd ed.,
Pearson, Boston, 2011.
\index{Artin, M.}

\bibitem{jacobson}
N.~Jacobson,
\textit{Basic Algebra I\,\&\,II}, 2nd ed.,
W.~H.~Freeman, New York, 1985--1989.
\index{Jacobson}

\bibitem{bourbaki}
N.~Bourbaki,
\textit{Alg\`ebre}, Chapitres 4 \`a 7,
Springer-Verlag, Berlin, 2007.
\index{Bourbaki}

\bibitem{perrin}
D.~Perrin,
\textit{Alg\`ebre: Cours de math\'ematiques pures},
\'Editions de l'\'Ecole Polytechnique, 2019.
\index{Perrin}

\bibitem{stewart}
I.~Stewart,
\textit{Galois Theory}, 4th ed.,
CRC Press, Boca Raton, 2015.
\index{Stewart}

\end{thebibliography}


%% ============================================================
%%  INDEX AND END
%% ============================================================
\printindex

\end{document}
